\documentclass[11pt,a4paper]{ltjsarticle}

% 09_typed-fcp-fkp 日本語版（査読用の全訳）
% 正本は英語版 ../main.tex（Discussion はコミット 8c7bd628 に対応）。
% 定理・補題・命題などの番号、ラベル、数式は英語版と一対一に対応させてある。
% 英語版と対応する。序論の冒頭と例 1.1 は著者による日本語版の改稿を英語版に反映したもの。

\usepackage{lmodern}
\usepackage{amsmath,amssymb,amsthm,mathtools}
\numberwithin{equation}{section}
\usepackage{enumitem}
\usepackage[margin=2.5cm]{geometry}
\usepackage{booktabs}
\usepackage{array}
\usepackage[round,authoryear]{natbib}
\usepackage[hidelinks]{hyperref}
\usepackage{url}
\usepackage[switch]{lineno}
\setlength{\emergencystretch}{3em}

\theoremstyle{plain}
\newtheorem{theorem}{定理}[section]
\newtheorem{lemma}[theorem]{補題}
\newtheorem{proposition}[theorem]{命題}
\newtheorem{corollary}[theorem]{系}
\newtheorem{fact}[theorem]{事実}
\theoremstyle{definition}
\newtheorem{definition}[theorem]{定義}
\newtheorem{example}[theorem]{例}
\theoremstyle{remark}
\newtheorem{remark}[theorem]{注意}
\newtheorem{openproblem}[theorem]{未解決問題}
\theoremstyle{plain}
\newtheorem{mainthm}{定理}
\renewcommand{\themainthm}{\Alph{mainthm}}
\renewcommand{\proofname}{証明}

\newcommand{\eps}{\lambda}
\newcommand{\R}{^{\rhd}}
\newcommand{\Lc}{^{\lhd}}
\newcommand{\RL}{^{\rhd\lhd}}
\newcommand{\cl}{\operatorname{cl}}
\newcommand{\D}{\mathcal D}
\newcommand{\Sub}{\operatorname{Sub}}
\newcommand{\Con}{\operatorname{Con}}
\newcommand{\Ccf}[1]{\mathcal C^{\mathrm{cf}}_{#1}}
\newcommand{\one}{\mathbf 1_\sharp}
\newcommand{\Fib}[2][]{F^{#1}_{#2}}
\newcommand{\FCP}[2]{\mathrm{FCP}^{#1}_{\mathrm{#2}}}
\newcommand{\FKP}[2]{\mathrm{FKP}^{#1}_{\mathrm{#2}}}
\newcommand{\tFCP}[3]{\mathrm{FCP}^{#1}_{\mathrm{#2}}[#3]}
\newcommand{\tFKP}[3]{\mathrm{FKP}^{#1}_{\mathrm{#2}}[#3]}
\newcommand{\Lfl}{\lfloor}
\newcommand{\Rfl}{\rfloor}
\newcommand{\Lce}{\lceil}
\newcommand{\Rce}{\rceil}
\newcommand{\ta}{\mathtt a}

\title{枠 vs. 見本：統語カテゴリーの指数的な隔たり}
\author{Takayuki Kuriyama\\
独立研究者（東京，日本）\\
\texttt{growup.kuriyama@gmail.com}}
\date{}

\begin{document}
\linenumbers
\maketitle

\begin{abstract}
統語論への分布的なアプローチでは，カテゴリーを有限の方法で名指す仕方が二つある．
そのメンバーが現れる少数の文の枠によって名指すか，少数の見本によって名指すかである．
後者では，見本に共通する文脈を考え，それらすべての文脈に出現可能な表現の集合をカテゴリーとする．
Kanazawa and Yoshinaka (2017) は，対応する学習器が成功するための有限文脈性と有限核性を特定した．
本稿では，この二つの仕方が指数的に異なりうることを示す．
二項の結合と一項の標識付けで作られ，どの枝もちょうど $n$ 段である構造について，標識が少なくとも一つあることを要求する．
非終端記号の数と，右辺と証人の長さで測ると，ある文法はその $n+1$ 個のカテゴリーをそれぞれ長さ $2n$ 以下の一つの枠で名指し，右辺の長さは $3n$ 以下だが生成規則は指数個である．
これに対し，どの見本の提示でもこの三つの積は $n$ について指数的であり，$O(n)$ 倍の範囲で最適である．このことから特定の学習器のデータに指数的な下界が得られる．
すべての生成規則と証人を数える全長で測ると，表現が標識を持つかどうかを記録する固定された観測器のもとでは全長 $O(n^2)$ の枠の提示があるのに対し，
どの観測器のもとでも見本の提示は指数的な全長を持つ．$\Lce\ta\Rce$ と $\Lfl\ta\ta\Rfl$ を分ける観測器がなければ，枠の提示も指数的な全長を要する．
二つの界は一般の一語の削除に拡がり，下界は接頭辞の長さで表される．
観測器を細分化すると，有限文脈性と有限核性のすべてではないがいくつかについて，枠には有利にしか働かず，見本には不利にしか働かない．
\end{abstract}

\noindent\textbf{キーワード：} 分布学習；統語概念束；有限文脈性；有限核性；
有限半群；特性標本．

\section{序論}\label{sec:intro}
\paragraph{分布からのカテゴリー.}
\citet{harris1954} の分布主義の伝統では，統語カテゴリーとは共通の文脈に出現できる表現の集合である．
構造主義の文法は少数の試験枠への代入によって語類を定め \citep[pp.~74--75]{fries1952}，事例基盤のモデルは蓄えられた事例によって働き \citep{bod2006}，用例基盤のモデルはそこからカテゴリーが生じるとする \citep{bybee2006}．
\citet{clark2009fg,clark2010icgi,clark2013} の統語概念束はこの考え方を定式化したものである．
言語 $L$ を固定する．\emph{枠}または\emph{文脈}とは文字列の対 $(u,v)$ であり，表現 $x$ がそれに入るとは $uxv\in L$ のことである．
枠の集合 $C$ に対し $C\Lc$ はすべての枠に入る表現の集合，表現の集合 $K$ に対し $K\R$ はすべての表現が入る枠の集合である．
したがってカテゴリーは，有限個の\emph{枠}によって $C\Lc$ として名指すことも，有限個の\emph{見本}によって $K\RL$，
すなわち見本に共通するすべての枠に入る表現の集合として名指すこともできる．
文法 $G$ が\emph{有限文脈性}（FCP）を持つとは有限集合 $C_A$ で $L(G,A)=C_A\Lc$ となること，
\emph{有限核性}（FKP）を持つとは有限集合 $K_A$ で $L(G,A)=K_A\RL$ となることである
\citep{clark2010icgi,yoshinaka2011dlt,clark-yoshinaka2016}．
\citet{ky2017} は，その\emph{極弱}の版（第~\ref{sec:prelim} 節）が，所属質問を用いる枠に基づく（双対）学習器と見本に基づく（主）学習器の成功を特徴づけることを示した．
本稿の見本はこれらの学習器の核であって，用例基盤の文献の類似度に基づく事例ではない．二つの伝統に共通するのは，カテゴリーをその環境で名指すか，その事例で名指すかという選択である．

\begin{example}[数の一致を持つ小さな断片]\label{ex:toy}
単語 \emph{the, dog, dogs, John, barks, bark, saw} 上の言語 $L_{\mathrm{ag}}$ を次の文法で生成されるものとする．
\begin{gather*}
S\to NP_{\mathrm{sg}}\ \textit{barks}\mid
NP_{\mathrm{pl}}\ \textit{bark}\mid NP\ \textit{saw}\ NP,\\
NP\to NP_{\mathrm{sg}}\mid NP_{\mathrm{pl}},\\
NP_{\mathrm{sg}}\to\textit{the dog}\mid\textit{John},\qquad
NP_{\mathrm{pl}}\to\textit{the dogs}.
\end{gather*}
単数名詞句は一つの枠でも一つの見本でも名指される：
$\{(\eps,\textit{barks})\}\Lc=\{\textit{John}\}\RL=\{\textit{the dog},\textit{John}\}$．
名詞句全体は一つの枠で名指されるが，見本が一つでは数の一致の制約が残り，単数と複数の見本が一つずつ要る：
$\{(\textit{John saw},\eps)\}\Lc=\{\textit{the dog},\textit{the dogs}\}\RL=\{\textit{the dog},\textit{John},\textit{the dogs}\}$．
\end{example}

\paragraph{問いと答え.}
文脈自由言語の有限の提示は，カテゴリーを枠で名指すときと見本で名指すときとで，それぞれどれだけ大きくなければならないか．
提示の大きさは，非終端記号の数 $|N|$，右辺の最大長 $\rho$，見本の最大長 $\kappa$ または枠の最大長 $\ell=|u|+|v|$ で測る．これらは学習器のデータの下界が通る量である（第~\ref{sec:applications} 節）．
別に，生成規則も数える全長でも測る（定義~\ref{def:presentation}）．

構造のどこかにある一つの構成素が標識を持てばその構造を認める\emph{大域的な認可条件}を考える．
二項の結合を $\Lfl xy\Rfl$，標識付きの構成素を $\Lce x\Rce$ と書き，
根から葉までのどの道もちょうど $n$ 個の構成子を通り，\emph{少なくとも一つの標識}を含む構造の集合を $M_n$ とする．
この存在の形の条件はよく知られている．語の韻律では，\emph{義務性}がどの語彙語にも第一強勢を持つ音節を少なくとも一つ要求し，強勢アクセント体系を定義する性質とされる \citep{hyman2006}．
厳格層仮説は各韻律水準をすぐ下の水準の構成素から組み立てるので （\citealp[p.~26]{selkirk1984}；\citealp[p.~7]{nespor-vogel1986}），固定した二つの水準のあいだのどの経路も同じ数の層を通る．
標準日本語では，節は一般に少なくとも一つの主格項を含むことを求められ，斜格主語構文のような例外がある \citep{kishimoto2010}．また複数の主格も許される \citep[Ch.~3]{kuno1973}．
葉の記号が一種類で深さが一様なら，標識のない構造はちょうど一つなので，$M_n$ はそのような構造すべてからなる言語 $L_n$ から一語を除いたものである：この削除は認可条件そのものであり，恣意的な例外ではない．
本稿はこの種の条件の論理的な形を取り出す．ここでの標識は一項の構成子であって卓立のラベルではなく，強勢も格もモデル化しない．
標識付きの構成素が一つあれば所属には十分だが，標識のない構成素は，認可される構造にもただ一つ除外される構造にも現れる．
枠はこれを容易に扱える：文法は最上位の標識の位置ごとに一つずつ根の生成規則を置いて標識を導入し，ほかのカテゴリーには制限を加えなくてよい．
見本も同じカテゴリーを名指せるが，$w^*_{n-1}$ ほど長い語を要する．

\begin{mainthm}[枠と見本；命題~\ref{prop:deletion-upper}，\ref{prop:untyped-frames}，定理~\ref{thm:deletion}]
$n\ge2$ とする．
\begin{enumerate}[label=\textup{(\roman*)},nosep]
\item 非終端記号が $n+1$ 個，右辺の長さが $3n$ 以下，生成規則が $2^n+2n-2$ 個の $M_n$ の文法 $F_n$ があって，そのどのカテゴリーも長さ $2n$ 以下の一つの枠で名指される．
\item 極弱 FKP を持つ $M_n$ のどの CFG $G$ も $\rho_G\,|N_G|\,\max(\kappa_G,1)\ge3\cdot2^n-n-5$ を満たす．
\item 非終端記号が $2n$ 個で右辺の長さが $4$ 以下の $M_n$ の文法 $H_n$ があって，
非終端記号あたり一つの，長さ $\max(3\cdot2^{n-1}-2,\,2n+1)$ 以下の見本で極弱 FKP を持つ．
\end{enumerate}
\end{mainthm}

削除がなければ隔たりはない：そのような構造すべてからなる言語 $L_n$ は，非終端記号 $O(n)$ 個，右辺の長さ $4$ 以下，証人の長さ $O(n)$ の枠と見本の提示を持つ（命題~\ref{prop:deletion-upper}(a)）．
理由は\emph{捕獲}である（補題~\ref{lem:capture}）：補題~\ref{lem:one-word} の標識付きの語 $RT$（$R$ は $w^*_n$ の標識のない接頭辞）で，$R$ の内側にある節点の産出はどれも，その節点のカテゴリーの見本でなければならないが，枠はそれ自身が標識を含みうる．定理~A は生成規則を数えない．全長での隔たりは定理~C である．

\paragraph{有限観測.}
表現そのものについての有限の情報も用いる．これを有限半群への全射準同型 $s:\Sigma^+\to\mathsf S$ でモデル化し，
\emph{型付け}または\emph{観測器}と呼ぶ．$s(w)$ は $w$ の\emph{型}であり，複合表現の型は構成要素の型から決まる．
\citet[Def.~2.1]{kuriyama-tcs1}%
\footnote{\citet{kuriyama-tcs1} の結果の番号は，その arXiv 第 5 版（arXiv:1409.6247v5）による．}
にならい，$x,y\in\Sigma^+$ について $s(x)=s(y)$ かつ $\{x\}\R\cap\{y\}\R\ne\emptyset$ ならば $\{x\}\R=\{y\}\R$ となるとき，
$L$ は \emph{$\sim_s$-代入可能}であるといい，文脈自由で $\sim_s$-代入可能な言語のクラスを $\Ccf s$ と書く．
自明な型付け $s_0$ のとき，これは \citet{clark-eyraud2007} の代入可能性を非空の断片に制限したものである．
例~\ref{ex:toy} で，$w$ が \emph{dogs} または \emph{bark} を含むとき $s(w)=\zeta$，そうでないとき $s(w)=1$ とする（$U_1=\{1,\zeta\}$，$\zeta$ は吸収元）．
型付きの枠はカテゴリーを一つの型に制限できる：$\{(\textit{John saw},\eps)\}\Lc\cap s^{-1}(1)$ は単数名詞句の集合であり，一致の枠なしに名指される．
一方，型付きの見本はすべてそのカテゴリーの型を持たなければならないので，\emph{the dog} と \emph{the dogs} を合わせて名詞句全体を名指すことはもうできない．

\paragraph{学習.}
\citet{kuriyama-tcs1} の学習モデルは正例だけを用いる．所属質問は言語学の言葉でいえば文法性判断の要求であり，
子どもの言語獲得では乏しいと長く考えられてきた種類の否定証拠である \citep{marcus1993}．
各クラス $\Ccf s$ は正例から極限同定できるが，その和はすべての正則言語を含む \citep[Prop.~3.1]{kuriyama-tcs1} ので同定できない \citep{gold1967}．
所属質問があれば，型なしの主学習器一つが和の全体を同定する（注意~\ref{rem:uniform-known}）．
定理~A は見本でカテゴリーを名指す学習器を下から抑える（第~\ref{sec:applications} 節）：$M_n$ 上で，\citet{kuriyama-tcs1} の再構成 $\mathsf{Rec}_s$ は
$M_n\in\Ccf s$ となるどの型付け $s$ についても大きさ $(3\cdot2^n-n-5)^{1/3}$ 以上の特性データを要し，パラメータを固定した \citet{ky2017} の主学習器は $2^{\Omega(n)}$ を要し，
別の論法により，双対学習器（型なしまたは $s_U$，階数 $r\ge2$）は収束時に長さ $3\cdot2^{n-1}-2$ 以上のデータを要する（系~\ref{cor:rec-data}，~\ref{cor:primal-data}，命題~\ref{prop:dual-refutation}）．

\paragraph{観測器の細分化.}
$h(x)=h(y)$ ならば $s(x)=s(y)$ となるとき，型付け $h$ は $s$ を\emph{細分化する}といい，$s\preceq h$ と書く．
最も粗い型付けは $s_0$ である．\emph{2 元の指示型付け}は，上の複数形の観測器のように，表現が与えられた集合の文字を含むかどうかだけを記録する．
極弱 $k$-FCP，$k$-FKP を持つ CFG を持つ言語のクラスを $\FCP{k}{vw}$，$\FKP{k}{vw}$ と書き，$\mathrm{st}$ と $\mathrm{w}$ で強い版と弱い版を，$[s]$ を添えて，各カテゴリーが型を持ちその型に制限される型付き版を表す（第~\ref{sec:prelim} 節と第~\ref{sec:typed} 節）．

\begin{mainthm}[文脈と核の双対性；命題~\ref{prop:duality}
と~\ref{prop:strict-context}--\ref{prop:weak-not-preserved}]
$s\preceq h$ かつ $k\ge1$ とする．このとき
\[
\tFKP{k}{vw}{h}\subseteq\tFKP{k}{vw}{s}\subseteq\FKP{k}{vw},\qquad
\FCP{k}{vw}\subseteq\tFCP{k}{vw}{s}\subseteq\tFCP{k}{vw}{h}
\]
であり，強 FCP についても同じ包含が成り立ち，型付きの極弱 FKP と弱 FKP は一致する．
すべての $k\ge1$ について，自明な型付けと 2 元の指示型付けのあいだの型付きの包含は，両側とも真である．
弱 FCP は単調でない．ある言語は，型なしでも自明な型付けについても弱 $2$-FCP を持つが，
2 元の指示型付けについては型付き弱 FCP を持たない．
\end{mainthm}

\paragraph{標識の観測器が加えるもの.}
表現が標識を含むかどうかを記録する 2 元の観測器を $s_U$ とする．

\begin{mainthm}[標識の観測器；第~\ref{sec:duality} 節]
$n\ge2$ とする．
\begin{enumerate}[label=\textup{(\roman*)},nosep]
\item $s_U$ のもとで，定理~A の文法 $H_n$ のどのカテゴリーも長さ $2n$ 以下の一つの枠で名指される．
$H_n$ の生成規則は $6n-4$ 個であり，$F_n$ では $2^n+2n-2$ 個である．
\item どの観測器のもとでも，$M_n$ のどの型付きの見本の提示も定理~A(ii) の下界を満たす．
$s_U$ のもとでは $L_n$ でさえ長さ $3\cdot2^n-2$ の見本を要する．観測器がなければ $2n+1$ で足りる．
\item 全長で測ると，(i) の枠の提示は $O(n^2)$ であり，$M_n$ のどの見本の提示も，どの観測器で型付けしても型なしでも，
$2\sqrt{3\cdot2^n-n-5}$ 以上である．
\item 観測器がないとき，および $\Lce\ta\Rce$ と $\Lfl\ta\ta\Rfl$ に同じ型を与えるどの観測器のもとでも，$M_n$ のどの枠の提示も全長 $2^{n/2}$ 以上を持つ．
\end{enumerate}
\end{mainthm}

(iii) と (iv) により，観測器は全長での隔たりを生み，それに必要である．観測器の選び方は重要である：長さの偶奇の観測器は $L_n$ の見本に不利に働かない（注意~\ref{rem:Ln-sU}）．
表~\ref{tab:summary} に評価をまとめる．

\begin{table}[t]
\centering\small\renewcommand{\arraystretch}{1.15}
\caption{$L_n$ と $M_n$ の評価．枠の提示の大きさは $\rho|N|\ell$，見本の提示の大きさは $\rho|N|\max(\kappa,1)$，全長は定義~\ref{def:presentation} による．下界はすべての提示について，上界は挙げた文法について成り立つ．}
\label{tab:summary}
\begin{tabular}{@{}>{\raggedright\arraybackslash}p{3.1cm}>{\raggedright\arraybackslash}p{3.6cm}>{\raggedright\arraybackslash}p{8.6cm}@{}}
\toprule
 & $L_n$ & $M_n=L_n\setminus\{w^*_n\}$\\
\midrule
枠，観測器なし & $O(n^2)$（$G_n$） & $O(n^3)$（$F_n$，命題~\ref{prop:untyped-frames}）；全長 $\ge2^{n/2}$（定理~\ref{thm:frame-total}），$F_n$ は $\Theta(n2^n)$\\
枠，$s_U$ & $O(n^2)$（注意~\ref{rem:Ln-sU}） & 全長 $O(n^2)$（$H_n$，命題~\ref{prop:mark-total}）\\
見本，観測器なし & $O(n^2)$（$G_n$） & $\ge3\cdot2^n-n-5$（定理~\ref{thm:deletion}），$H_n$ は $\le8n\max(3\cdot2^{n-1}-2,2n+1)$；全長 $\ge2\sqrt{3\cdot2^n-n-5}$（命題~\ref{prop:mark-total}）\\
見本，任意の観測器 & $s_U$ のもとで $\kappa\ge3\cdot2^n-2$（注意~\ref{rem:Ln-sU}） & 同じ下界\\
$\mathsf{Rec}_s$，正例 & $s_0$ で多項式，$s_U$ で指数的 \citep{kuriyama-tcs1} & $M_n\in\Ccf s$ なるすべての $s$ について $\|K\|^3\ge3\cdot2^n-n-5$（系~\ref{cor:rec-data}）\\
$\mathrm{Primal}_{r,k}$，収束時 & --- & $\|D_i\|\ge((3\cdot2^n-n-5)/(2r+1))^{1/(2k+2)}$（系~\ref{cor:primal-data}）\\
双対学習器（$r\ge2$，型なしまたは $s_U$），収束時 & --- & 長さ $3\cdot2^{n-1}-2$ 以上のデータ（命題~\ref{prop:dual-refutation}）\\
\bottomrule
\end{tabular}
\end{table}

枠の上界と見本の下界は一般に成り立つ．
枠は，初期記号がどの右辺にも現れず，初期でないすべてのカテゴリーの枠が型のイデアルを強制するとき，そのイデアルとの交わりにそのまま移り（命題~\ref{prop:ideal}），
有限言語 $L$ と $w=RT^\circ\in L$ について，ある $T\ne T^\circ$ で $RT\in L$ かつすべての $y$ について $yT\in L$ ならば $yT^\circ\in L$ となるなら，$L\setminus\{w\}$ のどの見本の提示も $\rho|N|\max(\kappa,1)\ge|R|$ を満たす
（補題~\ref{lem:one-word}，定理~\ref{thm:general-lower}）．
定理~\ref{thm:general-gap} は二つを合わせたものであり，$M_n$ のほかに，木構造を持たない平らな族もその例である（例~\ref{ex:flat}）．
定理~\ref{thm:frame-total} の方法はその族にも当てはまる．

\paragraph{既知のことと新しいこと.}
型付きの学習器は \citet{yoshinaka2015} のソート付き枠組みの例であり，本稿は用いる健全性だけを証明する（命題~\ref{prop:learners}）．
型付きの文脈集合は \citet{ky2023} の拡張正則閉包に入り，
\citet[Sec.~8]{yoshinaka2015} は異なるシグネチャについて学習可能なクラスをすでに比較している．
定理~B の包含は閉包の計算の短い帰結であり（補題~\ref{lem:star} と~\ref{lem:reduction}），
本稿が加えるのは枠と見本の逆向きの振る舞い，すべての $k$ についての真包含の証人，および弱 FCP の単調性の破れである．
族 $L_n$，その文法，見本 $u_i$ は \citet[Sec.~6]{kuriyama-tcs1} のものであり，そこでは $M_n$ は $L_n^-$ と書かれ，区別のための対の第二の成員として現れるだけである．
標識の観測器のもとでの $L_n$ の評価（注意~\ref{rem:Ln-sU}）は，\citet[式~($\spadesuit$)，Cor.~6.6]{kuriyama-tcs1} の評価を収量から核に移したものである．
隔たり，文法 $F_n$，捕獲補題と混合の補題，そこから得られる下界は，調べた限りでは新しい．\citet{ky2023} は枠の側だけを扱う．二つの名指し方のこれまでの比較は定性的であるか \citep{yoshinaka2011dlt,yoshinaka2012lata,ky2021}，ランダムな文法で各性質が成り立つ頻度を数えるもの \citep{clark2016icgi} であり，一つの言語の二つの提示の大きさを比べたものはない．
指数的な特性データそれ自体は新しくない \citep[Example~2]{clark-eyraud2007}（\citealp{coste-nicolas2019} も参照）．
しかしここでは長い語だけでは下界は生じない：$L_n$ は同じ長い語を含み，非終端記号 $O(n)$ 個と長さ $O(n)$ の見本からなる見本の提示を持つ．

%======================================================================
\section{準備}\label{sec:prelim}

\subsection*{文脈，閉包，健全な前不動点}
$\Sigma$ を有限アルファベット，$\eps$ を空語とし，$\Sigma^+=\Sigma^*\setminus\{\eps\}$ とする．$L\subseteq\Sigma^*$ を固定する．
$C\subseteq\Sigma^*\times\Sigma^*$ と $K\subseteq\Sigma^*$ に対して
\[
C\Lc=\{x\in\Sigma^*: uxv\in L\ \forall(u,v)\in C\},\qquad
K\R=\{(u,v): uKv\subseteq L\}
\]
とおく．したがって $\emptyset\Lc=\Sigma^*$，$\emptyset\R=\Sigma^*\times\Sigma^*$ である．
$\D_L(x)=\{x\}\R$ と書く．対 $(\cdot\R,\cdot\Lc)$ はガロア接続である：$K\subseteq C\Lc\iff K\R\supseteq C$．
特に，つねに $K\subseteq K\RL$ かつ $C\subseteq C\Lc{}\R$ である．
よって $K\mapsto K\RL$ は閉包作用素であり，$K\R=K\RL{}\R$，かつ $Y\subseteq K\RL\iff Y\R\supseteq K\R$ である．
これらの作用素は $L$ に依存する．複数の言語が現れるときは「$L$ に相対的に」という．

\begin{lemma}[Clark の包含；\citealp{clark2009fg} を参照]\label{lem:clark}
すべての $X,Y\subseteq\Sigma^*$ について $X\RL Y\RL\subseteq(XY)\RL$．
したがって，終端語 $w_0,\dots,w_n$ と集合 $X_1,\dots,X_n$ について
$w_0X_1\RL w_1\cdots X_n\RL w_n\subseteq(w_0X_1w_1\cdots X_nw_n)\RL$．
\end{lemma}
\begin{proof}
$x\in X\RL$，$y\in Y\RL$，$(u,v)\in(XY)\R$ なら，すべての $x'\in X$ について $(ux',v)\in Y\R$ なので $(u,yv)\in X\R$，よって $uxyv\in L$．
第二の主張は帰納法による．
\end{proof}

\emph{文脈自由文法}（CFG）とは，初期非終端記号の集合 $I\subseteq N$ を持つ $G=(N,\Sigma,P,I)$ のことである \citep{ky2017}．
$A$ の終端収量を $L(G,A)$，$L(G)=\bigcup_{S\in I}L(G,S)$，
$\mathcal C(G,A)=\{(u,v): \text{ある } S\in I \text{ について } S\Rightarrow^*_G uAv\}$ と書く．

族 $(X_A)_{A\in N}$ が $G$ の\emph{前不動点}であるとは，すべての生成規則 $A\to w_0A_1w_1\cdots A_nw_n$ について
$X_A\supseteq w_0X_{A_1}w_1\cdots X_{A_n}w_n$ が成り立つことをいう．
族 $(L(G,A))_A$ は最小の前不動点であるから，任意の前不動点について $L(G,A)\subseteq X_A$ である．
さらに $\bigcup_{S\in I}X_S=L(G)$ が成り立つとき，\emph{健全な前不動点}（SPP）という．
すなわち，健全な前不動点は，生成規則との整合性を保ちつつ，文法全体の生成言語を変えない範囲で，
各非終端記号の生成言語の過大評価を許す．

例えば $S\to bAc\mid baaac$，$A\to a\mid aa$ について，族 $X_A=\{a,aa,aaa\}$，$X_S=L(G,S)=\{bac,baac,baaac\}$ は
$L(G,A)\subsetneq X_A$ なる SPP である：余分な語 $aaa$ が生むのは，いずれにせよ生成される $baaac$ だけである．
$S\to baaac$ を除くと，この族は前不動点のままだが健全ではない．

\begin{lemma}[到達文脈]\label{lem:reach}
$(X_A)_A$ が $G$ の SPP で $(u,v)\in\mathcal C(G,A)$ ならば，$uX_Av\subseteq L(G)$．
\end{lemma}
\begin{proof}
文形式の各非終端記号 $B$ を $X_B$ で置き換えたものは，$(X_B)$ が前不動点なので導出の一歩ごとに縮むだけである．
よって $S\Rightarrow^*uAv$ から $uX_Av\subseteq X_S\subseteq L(G)$．
\end{proof}

\citet[Sec.~3]{ky2017} にならい，$L=L(G)$ に相対的に，集合 $X$ が\emph{$k$-文脈生成}であるとは $|C|\le k$ なる $C$ で $X=C\Lc$ となることをいい，
\emph{$k$-核生成}であるとは $|K|\le k$ なる $K$ で $X=K\RL$ となることをいう．
文法 $G$ が\emph{強}（\emph{弱}）$k$-FCP を持つとは，$(L(G,A))_A$（それぞれ $(L(G,A)\RL)_A$）が $k$-文脈生成であることをいい，
\emph{極弱} $k$-FCP を持つとは，$G$ が $k$-文脈生成の SPP を持つことをいう．
FKP についても「文脈」を「核」に替えて同様に定める．
強，弱，極弱の性質を持つ CFG がある言語のクラスを，それぞれ $\FCP{k}{st}$，$\FCP{k}{w}$，$\FCP{k}{vw}$
（核についても同様に $\FKP{k}{st}$，$\FKP{k}{w}$，$\FKP{k}{vw}$）と書く．
$k$ を省いたときは，すべての $k\ge1$ にわたる和を表す．
$(C\Lc)\RL=C\Lc$，$(K\RL)\RL=K\RL$ なので強い性質は弱い性質を導き，
$(L(G,A)\RL)_A$ はつねに $G$ の SPP なので弱い性質は極弱の性質を導く
（補題~\ref{lem:clark} による．健全性は，初期記号 $A$ について $L(G,A)\subseteq L$ から $(\eps,\eps)\in L(G,A)\R$ となることによる）．
したがって
\[
\FCP{k}{st}\subseteq\FCP{k}{w}\subseteq\FCP{k}{vw},\qquad
\FKP{k}{st}\subseteq\FKP{k}{w}\subseteq\FKP{k}{vw}
\]
であり，二つ目の包含は核については等号，文脈については $k\ge2$ で真の包含である（事実~\ref{fact:ky}(a)(c)）．

\begin{fact}[\citealp{ky2017}]\label{fact:ky}
(a)（Prop.~10）$\FKP{k}{vw}=\FKP{k}{w}$．より正確には，$(K_A\RL)_A$ が $G$ の SPP ならば，
生成規則 $A\to w$ $(w\in K_A)$ を加えると，$L(G')=L(G)$ かつ $L(G',A)\RL=K_A\RL$ となる $G'$ が得られる．
(b)（Theorems~3--6）パラメータ $r,k$ の双対（主）学習器は，$L_*$ の正の提示と所属オラクルを与えられたとき，
$L_*$ が右辺あたり高々 $r$ 個の非終端記号を持ち極弱 $k$-FCP（$k$-FKP）を持つ CFG を持つならば，$L_*$ の文法に収束する．
また，収束先のどの文法も極弱 $k$-FCP（$k$-FKP）を持ち，その SPP は自身の非終端記号で添字づけられた $(C\Lc)_C$
（それぞれ $(K\RL)_K$）である．
(c)（Prop.~2，Thm.~14）$\FKP1{st}\setminus\FCP{}{vw}$，$\FCP1{st}\setminus\FKP{}{vw}$，$\FCP2{vw}\setminus\FCP{}{w}$
のそれぞれに言語が存在する．
\end{fact}

\paragraph{学習器.}
\citet{ky2017} の Algorithms~1 と~2，すなわち\emph{双対学習器} $\mathrm{Dual}_{r,k}$ と\emph{主学習器} $\mathrm{Primal}_{r,k}$（$r,k\ge1$）を思い出す．
どちらも標的 $L_*$ の正の提示 $t_1,t_2,\dots$ を読み，所属質問を行い（$\Lc$ と $\R$ は $L_*$ に相対的），各段階 $i$ で文法 $G_i$ を出力する．
材料は $D_i=\{t_1,\dots,t_i\}$，$D$ の語の部分文字列の集合 $\Sub(D)$，$\Con(D)=\{(u,v):\text{ある }w\text{ について }uwv\in D\}$，
およびある $u_i$ について $u_0w_1u_1\cdots w_lu_l\in D$ となる組 $(w_1,\dots,w_l)$，$1\le l\le m$ の集合 $\Sub^{\le m}(D)$ である．
$D_i\not\subseteq L(G_{i-1})$ のとき段階 $i$ を\emph{リセット}と呼ぶ．
リセットでは $\mathcal E_i=\Sub(D_i)$，$J_i=\Con(D_i)$，$H_i=\Sub^{\le r+1}(D_i)$ を計算し直し，それ以外では前の値を保つ．
ただし双対学習器は $\mathcal E_i$ を，主学習器は $J_i$ を毎段階計算し直す．
生成規則は $(w_0,\dots,w_n)\in H_i$ なる $A_0\to w_0A_1w_1\cdots A_nw_n$ の形である．
双対学習器の非終端記号は $1\le|C|\le k$ なる $C\subseteq J_i$，生成規則は\emph{$\mathcal E_i$ 上で妥当な}もの，すなわち
$C_0\Lc\supseteq w_0(\mathcal E_i\cap C_1\Lc)w_1\cdots(\mathcal E_i\cap C_n\Lc)w_n$ を満たすもの，初期非終端記号は $\mathcal E_i\cap C\Lc\subseteq L_*$ となる $C$ である．
主学習器の非終端記号は $1\le|K|\le k$ なる $K\subseteq\mathcal E_i$，生成規則は $J_i\cap K_0\R\subseteq(w_0K_1w_1\cdots K_nw_n)\R$ を満たすもの，初期非終端記号は $K\subseteq L_*$ である．
すべての $j\ge i$ について $G_j=G_i$ となるとき，学習器は\emph{段階 $i$ で収束する}という．
（\citet{ky2017} の $E_i$ を $\mathcal E_i$ と書く．）

\subsection*{固定された有限半群の型付け}
\emph{型付け}または\emph{観測器}とは，有限半群への全射な半群の射 $s:\Sigma^+\to\mathsf S$ のことである．
新しい単位元を加えた $\mathsf S^\sharp=\mathsf S\cup\{\one\}$ と，そのモノイドとしての拡張 $\widehat s:\Sigma^*\to\mathsf S^\sharp$，
$\widehat s(\eps)=\one$ を考える \citep[Lemma~2.2]{kuriyama-tcs1}．
$\mu\in\mathsf S^\sharp$ について\emph{ファイバー}を $F_\mu=\widehat s^{-1}(\mu)$ とする．
したがって $F_{\one}=\{\eps\}$ であり，$\{F_\mu\}$ は $\Sigma^*$ を正則集合に分割する．
自明な型付けを $s_0:\Sigma^+\to\{0\}$ とする．

$h$ が $s$ を細分化するとき（$s\preceq h$，第~\ref{sec:intro} 節），$\widehat s=\pi\circ\widehat h$ となるモノイドの射 $\pi:\mathsf H^\sharp\to\mathsf S^\sharp$ がただ一つ存在し
（$\pi(h(w))=s(w)$，$\pi(\one)=\one$），$F^h_\nu\subseteq F^s_{\pi(\nu)}$ である．
文字の集合 $Z\subseteq\Sigma$，$\emptyset\ne Z\ne\Sigma$ について，\emph{$Z$-指示型付け} $h_Z:\Sigma^+\to U_1=\{1,\zeta\}$
（$\zeta$ は吸収元）は，$w$ が $Z$ の文字を含むとき，かつそのときに限り $w$ を $\zeta$ に送る．これは $s_0$ を細分化する．

\begin{definition}[型付き文法]\label{def:typed}
\emph{$s$-型付き CFG} とは，CFG と $\tau:N\to\mathsf S^\sharp$ の組であって，
すべての生成規則 $A\to w_0A_1w_1\cdots A_nw_n$ が
$\tau(A)=\widehat s(w_0)\tau(A_1)\widehat s(w_1)\cdots\tau(A_n)\widehat s(w_n)$ を満たすものをいう．
\end{definition}
導出についての帰納法により，すべての $A$ について $L(G,A)\subseteq F_{\tau(A)}$ である．
$s\preceq h$ で $(G,\tau)$ が $h$-型付きならば，$(G,\pi\circ\tau)$ は $s$-型付きである．

\begin{lemma}[文法の型付き細分化]\label{lem:product}
$G=(N,\Sigma,P,I)$ を CFG，$h$ を型付けとし，$\sigma:N\to\mathsf S^\sharp$ は未定義（型なしの場合）か，
$s\preceq h$ なる $s$ についての $G$ の $s$-型付けとする．
$G^h$ を次で定める．非終端記号は $\nu\in\mathsf H^\sharp$ についての $(A,\nu)$（型付きの場合は $\pi(\nu)=\sigma(A)$ のもの），
生成規則は $A\to w_0A_1w_1\cdots A_nw_n\in P$ と
$\nu=\widehat h(w_0)\nu_1\widehat h(w_1)\cdots\nu_n\widehat h(w_n)$ についての
$(A,\nu)\to w_0(A_1,\nu_1)w_1\cdots(A_n,\nu_n)w_n$，
初期非終端記号は $S\in I$ についての $(S,\nu)$，そして $\tau(A,\nu)=\nu$．
このとき $G^h$ は $h$-型付きであり，$L(G^h,(A,\nu))=L(G,A)\cap F^h_\nu$．特に $L(G^h)=L(G)$．
\end{lemma}
\begin{proof}
型の整合性は構成から成り立つ．包含 $\subseteq$ は $G^h$ での導出についての帰納法による．
$\supseteq$ は $G$ での導出についての帰納法により，収量 $x_i$ を持つ子に型 $\nu_i=\widehat h(x_i)$ を与えればよい
（型付きの場合は $\pi(\nu_i)=\widehat s(x_i)=\sigma(A_i)$）．
\end{proof}

%======================================================================
\section{型付き閉包と型付きの有限文脈性・有限核性}
\label{sec:typed}

$L$ と型付け $s$ を固定する．$\mu\in\mathsf S^\sharp$ と $K\subseteq F_\mu$ について，\emph{型付き閉包}を
\[
\cl_\mu(K)=K\RL\cap F_\mu
\]
で定める．

\begin{lemma}[型付きガロア接続]\label{lem:typed-galois}
$K\subseteq F_\mu$ と任意の $C$ について，$K\subseteq C\Lc\cap F_\mu\iff C\subseteq K\R$．
よって $\cl_\mu$ は $F_\mu$ の部分集合の上の閉包作用素であり，その閉集合はちょうど $C\Lc\cap F_\mu$ の形の集合であり，
$Y\subseteq F_\mu$ について $Y\subseteq\cl_\mu(K)\iff K\R\subseteq Y\R$ である．
\end{lemma}
\begin{proof}
同値は，型なしの同値を $K\subseteq F_\mu$ に制限したものである．
ほかの主張は，$\cl_\mu(K)=(K\R)\Lc\cap F_\mu$ と $\cl_\mu(K)\R=K\R$ から $\cdot\RL$ と同様に従う．
\end{proof}

次の等式は，\citet[Sec.~3, Eq.~(1)]{ky2017} の閉包と積の等式の直接の帰結を記録するものである．

\begin{lemma}[文脈集合は閉包に鈍感]\label{lem:star}
$w_0,\dots,w_n$ を終端語とし，$1\le i\le n$ について $X_i\subseteq Y_i\subseteq X_i\RL$ とする．このとき
\[
(w_0Y_1w_1\cdots Y_nw_n)\R=(w_0X_1w_1\cdots X_nw_n)\R .
\]
特に，$X_i\subseteq F_{\mu_i}$ のときの $Y_i=\cl_{\mu_i}(X_i)$ について，また $Y_i=X_i\RL$ について，これが成り立つ．
\end{lemma}
\begin{proof}
$P_X=w_0X_1\cdots X_nw_n$，$P_Y=w_0Y_1\cdots Y_nw_n$ とおく．$P_X\subseteq P_Y$ から $P_Y\R\subseteq P_X\R$．
補題~\ref{lem:clark} により $P_Y\subseteq w_0X_1\RL\cdots X_n\RL w_n\subseteq P_X\RL$ なので，
$P_Y\R\supseteq P_X\RL{}\R=P_X\R$．
\end{proof}

\begin{definition}[型付き FCP と FKP]\label{def:typed-props}
$(G,\tau)$ を $s$-型付きとし，$L(G)$ に相対的に考える．
$X_A\subseteq F_{\tau(A)}$ なる族 $(X_A)_A$ が\emph{型付き $k$-文脈生成}であるとは，$|C_A|\le k$ なる $C_A$ で $X_A=C_A\Lc\cap F_{\tau(A)}$ となることをいい，
\emph{型付き $k$-核生成}であるとは，$K_A\subseteq F_{\tau(A)}$，$|K_A|\le k$ なる $K_A$ で $X_A=\cl_{\tau(A)}(K_A)$ となることをいう
（核は\emph{型について均質}である）．
$(G,\tau)$ が型付きの強／弱／極弱 $k$-FCP（それぞれ $k$-FKP）を持つとは，
$(L(G,A))_A$／$(\cl_{\tau(A)}(L(G,A)))_A$／ある SPP が型付き $k$-文脈生成（それぞれ $k$-核生成）であることをいう．
対応する言語のクラス「$L$ はその性質を持つ $s$-型付き CFG を持つ」を
$\tFCP{k}{st}{s}$，$\tFCP{k}{w}{s}$，$\tFCP{k}{vw}{s}$，および $\tFKP{k}{\cdot}{s}$ で表す．
\end{definition}

\begin{definition}[提示とその大きさ]\label{def:presentation}
$L$ の \emph{$k$-文脈の提示}（\emph{枠の提示}）とは，$L(G)=L$ なる CFG $G$ と，
$(C_A\Lc)_A$ が $G$ の SPP となる $|C_A|\le k$ の有限集合 $C_A$ の組である．
\emph{$k$-核の提示}（\emph{見本の提示}）は，有限集合 $K_A$，$|K_A|\le k$ と SPP $(K_A\RL)_A$ について同様に定める．
\emph{型付きの}提示は，$s$-型付きの $(G,\tau)$ と，SPP $(C_A\Lc\cap F_{\tau(A)})_A$，
それぞれ $K_A\subseteq F_{\tau(A)}$ なる $(\cl_{\tau(A)}(K_A))_A$ を用いる．
SPP が $(L(G,A))_A$ であるとき，提示は\emph{強い}という．
提示の大きさは，非終端記号の数 $|N|$，右辺の記号の最大数 $\rho$，核の語の最大長 $\kappa$，
およびある $C_A$ に属する文脈 $(u,v)$ の長さ $|u|+|v|$ の最大値で測る．
提示の\emph{全長}は，すべての生成規則にわたる $\sum_{A\to\alpha}(|\alpha|+1)$ に，すべての核の語 $e$ にわたる $\sum(|e|+1)$，
それぞれすべての文脈 $(u,v)$ にわたる $\sum(|u|+|v|+1)$ を加えたものである．
和はすべての非終端記号 $A$ と $K_A$ または $C_A$ のすべての元にわたる．空集合の最大値は $0$ とし，
観測器，型 $\tau(A)$，初期記号の集合は数えない．
\end{definition}
したがって，$L$ が（$s$-型付きの）$k$-文脈の提示を持つのは $L\in\FCP{k}{vw}$（それぞれ $\tFCP{k}{vw}{s}$）のとき，かつそのときに限り，核についても同様である．

SPP をファイバーに制限しても SPP のままなので，型付きの SPP で失うものはない．
$s_0$-型付き文法は $\eps$ を非空の収量から分けるので型なしと完全に同じではない．型なしのクラスと比べるときはそれを明示する．

\begin{proposition}[型付き学習器の健全性]\label{prop:learners}
\emph{型付き双対学習器} $\mathrm{Dual}^s_{r,k}$ を，第~\ref{sec:prelim} 節の双対学習器で，
非終端記号を $(C,\mu)$，$C\subseteq J_i$，$1\le|C|\le k$，$\mu\in\mathsf S^\sharp$，
生成規則を型の整合するものだけ，妥当性の検査を
$C_0\Lc\cap F_{\mu_0}\supseteq w_0(\mathcal E_i\cap C_1\Lc\cap F_{\mu_1})w_1\cdots(\mathcal E_i\cap C_n\Lc\cap F_{\mu_n})w_n$，
初期非終端記号を $\mathcal E_i\cap C\Lc\cap F_\mu\subseteq L_*$ となる $(C,\mu)$ としたものと定める．
\emph{型付き主学習器} $\mathrm{Primal}^s_{r,k}$ を，主学習器で，非終端記号を $(K,\mu)$，$K\subseteq\mathcal E_i\cap F_\mu$，$1\le|K|\le k$，
生成規則を型の整合するものだけ，検査を型なしの妥当性の検査としたものと定める．
これらの学習器の一方が $L_*$ の正の提示の上で $G$ に収束するならば，$L(G)=L_*$ であり，
$(C\Lc\cap F_\mu)_{(C,\mu)}$（それぞれ $(\cl_\mu(K))_{(K,\mu)}$）は $G$ の SPP である．
主学習器については $(K\RL)_{(K,\mu)}$ も $G$ の SPP である．
\end{proposition}
\begin{proof}
\citet[Theorems~3 と~4 の証明]{ky2017} と同様である．段階 $i$ 以降の出力を $G$ とする．
$d\in L_*\setminus L(G)$ ならば，$d\in D_j$ なる段階 $j>i$ をとる．$D_j\not\subseteq L(G)=L(G_{j-1})$ なので段階 $j$ はリセットであり，$G_j$ は初期非終端記号 $(\{(\eps,\eps)\},\widehat s(d))$ と生成規則
$(\{(\eps,\eps)\},\widehat s(d))\to d$（双対），それぞれ $(\{d\},\widehat s(d))$ と $(\{d\},\widehat s(d))\to d$（主）を持ち，$G_j=G$ に反する．よって $L_*\subseteq L(G)$．
どの $C$ も $\Con(L_*)$ の空でない部分集合なので $C\Lc\subseteq\Sub(L_*)=\bigcup_j\mathcal E_j$ であり，どの $K$ も空でないので $K\R\subseteq\Con(L_*)=\bigcup_jJ_j$ である．
よってすべての段階 $j\ge i$ での検査により，この族は（主学習器については補題~\ref{lem:star} とガロア接続により）初期成分が $L_*$ に含まれる前不動点となり，
$L(G)\subseteq\bigcup_{S\in I}X_S\subseteq L_*\subseteq L(G)$．
主学習器については，補題~\ref{lem:reduction} が $(K\RL)$ を型付き SPP $(\cl_\mu(K))$ に変える．
\end{proof}
逆向き，すなわち \citet[Theorems~5--6]{ky2017} の型付きの類似は，本稿では証明も使用もしない．

次の補題は，隔たりと双対性の両方に用いられ，型付きの核の提示は核が型について均質な型なしの提示にほかならないことを述べる．

\begin{lemma}[核の縮約補題]\label{lem:reduction}
$(G,\tau)$ を $s$-型付き，$L=L(G)$ とし，各 $A\in N$ について $K_A\subseteq F_{\tau(A)}$ とする．
このとき，$(\cl_{\tau(A)}(K_A))_A$ が $G$ の SPP であるのは，$(K_A\RL)_A$ が $G$ の SPP であるとき，かつそのときに限る．
\end{lemma}
\begin{proof}
生成規則 $A\to w_0A_1w_1\cdots A_nw_n$ を固定し，$K_i=K_{A_i}$，
$P_0=w_0K_1w_1\cdots K_nw_n$，$P_1=w_0\cl(K_1)w_1\cdots\cl(K_n)w_n$，$P_2=w_0K_1\RL w_1\cdots K_n\RL w_n$ と書く．
補題~\ref{lem:star} により $P_0\R=P_1\R=P_2\R$．
$P_1\subseteq F_{\tau(A)}$（型の整合性）なので，補題~\ref{lem:typed-galois} により
$\cl_{\tau(A)}(K_A)\supseteq P_1\iff K_A\R\subseteq P_1\R$ であり，型なしのガロア接続により $K_A\RL\supseteq P_2\iff K_A\R\subseteq P_2\R$．
よって二つの前不動点の条件は，生成規則ごとに一致する．

健全性について．どちらの族も（一方がそうなら）前不動点であり，したがって $(L(G,A))_A$ を含むので，どちらの場合も $\bigcup_I X_S\supseteq L$．
$S\in I$ について，$K_S\subseteq L$ ならば $(\eps,\eps)\in K_S\R$ であり $\cl(K_S)\subseteq K_S\RL\subseteq L$．
逆に $K_S\subseteq\cl(K_S)$ かつ $K_S\subseteq K_S\RL$．
よってどちらの健全性の条件も，すべての $S\in I$ について $K_S\subseteq L$ となることと同値である．
\end{proof}
したがって，以下の見本の提示についての下界はすべて，どの型付けによる型付きの提示についても成り立つ：同じ文法の同じ核による型なしの SPP に移ればよい．

%======================================================================
\section{核と文脈の隔たり}\label{sec:deletion}

この節で定理~A を証明する．下界は一般の捕獲の条件のもとで証明する（定理~\ref{thm:general-lower}）．言語は有限なので，どの主張も生成能力ではなく提示の大きさについてのものである．

\citet[Sec.~6]{kuriyama-tcs1} の括弧付き倍化族を用いる．
$\Gamma=\{\ta,\Lfl,\Rfl,\Lce,\Rce\}$ とし，$G_n$ を
$S_0\to E_n$，$E_0\to\ta$，$E_i\to\Lfl E_{i-1}E_{i-1}\Rfl\mid\Lce E_{i-1}\Rce$ $(1\le i\le n)$ とし，$L_n=L(G_n)$ とする．
$w^*_0=\ta$，$w^*_i=\Lfl w^*_{i-1}w^*_{i-1}\Rfl$（したがって $|w^*_i|=3\cdot2^i-2$），$u_i=\Lce^i\ta\Rce^i$，
\[
M_n=L_n\setminus\{w^*_n\}
\]
とおく．$s_U:\Gamma^+\to U_1$ を $\{\Lce,\Rce\}$ の指示型付けとする．
このとき $w^*_i$ は $L(G_n,E_i)$ の型 $1$ のただ一つの語である：一項の節点がなければ，導出は高さ $i$ の完全二分木だからである．
定理~A とその証明はこれらの定義しか用いない．\citet{kuriyama-tcs1} による $M_n$ の学習可能性についての事実が必要になるのは第~\ref{sec:applications} 節だけである．

\begin{lemma}[節点の構造]\label{lem:nodes}
$m\in L_n$ とし，$y\in L(G_n,E_j)$ が $m$ に現れるとする．
このとき，その出現は $m$ の（ただ一つの）導出木の $E_j$-節点の収量であり，
それを任意の $y'\in L(G_n,E_j)$ で置き換えると $L_n$ の語が得られる．
\end{lemma}
\begin{proof}
$L_n$ の語は 2 種類の括弧について正しく括弧付けされており，対応する括弧の対はどれも $E_i$-節点の対であり，導出木は括弧付けから決まる．
$E_i$-節点の根から葉へのどの道もちょうど $i$ 個の括弧の対を含むので，$E_i$-節点の収量の入れ子の深さは $i$ である．
$j=0$ なら $y=\ta$ は葉である．
$j\ge1$ なら $y$ は釣り合っており，その最初の括弧は $y$ の中で最後の括弧と対応する（$y$ の真の非空接頭辞はどれも正の深さを持つ）．
したがって $m$ の中でも対応し，その出現は入れ子の深さ $j$ の節点，すなわち $E_j$-節点の収量である．
部分木を置き換えると $G_n$ の導出木が得られる．
\end{proof}

\begin{proposition}[上界]\label{prop:deletion-upper}
\begin{enumerate}[label=\textup{(\alph*)},nosep]
\item $L_n$ に相対的に，$L(G_n,E_i)=\{u_i\}\RL=\{(\Lce^{n-i},\Rce^{n-i})\}\Lc$．
よって $G_n$ は，長さ $\le2n+1$ の核と文脈で，強 $1$-FKP と強 $1$-FCP を持つ．
\item $M_n$ に相対的に，次の $s_U$-型付き文法 $H_n$ は $M_n$ を生成し，
文脈 $(\Lce^{n-i},\Rce^{n-i})$ $(i<n)$ と $(\eps,\eps)$ で型付き強 $1$-FCP を持つ．
非終端記号は $(E_i,1)$ $(0\le i\le n-1)$ と $(E_i,\zeta)$ $(1\le i\le n)$，
生成規則は $(E_0,1)\to\ta$，$(E_i,1)\to\Lfl(E_{i-1},1)(E_{i-1},1)\Rfl$ $(1\le i<n)$，
$(\alpha,\beta)\ne(1,1)$ についての $(E_i,\zeta)\to\Lfl(E_{i-1},\alpha)(E_{i-1},\beta)\Rfl$，および $(E_i,\zeta)\to\Lce(E_{i-1},\alpha)\Rce$，
初期非終端記号は $(E_n,\zeta)$ である．ここで $\zeta$ の規則は $1\le i\le n$ についてのもので，
子が上に示した非終端記号の集合に属する例だけを含める．非終端記号は $2n$ 個，生成規則は $6n-4$ 個，$\rho\le4$ である．
\item $M_n$ に相対的に，$H_n$ は $(E_i,1)$ には核 $\{w^*_i\}$，$(E_i,\zeta)$ には核 $\{u_i\}$ で，型付き強 $1$-FKP も持つ．
よって $(\{w^*_i\}\RL)$，$(\{u_i\}\RL)$ は $H_n$ の型なしの SPP であり，$H_n$ は
$\kappa=\max(|w^*_{n-1}|,|u_n|)=\max(3\cdot2^{n-1}-2,\,2n+1)$ で極弱 $1$-FKP を持つ．
\end{enumerate}
\end{proposition}
\begin{proof}
(a) $\Lce^{n-i}x\Rce^{n-i}\in L_n$ となるのは $x\in L(G_n,E_i)$ のとき，かつそのときに限る：
先頭の $\Lce$ が $E_n\to\Lce E_{n-1}\Rce$ を強い，以下同様である．
核について：$(u,v)\in\D(u_i)$ ならば，補題~\ref{lem:nodes} によりすべての $y\in L(G_n,E_i)$ について $uyv\in L_n$ なので，
$\D(y)\supseteq\D(u_i)$ であり，
$L(G_n,E_i)\subseteq\{u_i\}\RL\subseteq\{(\Lce^{n-i},\Rce^{n-i})\}\Lc=L(G_n,E_i)$．

(b) 帰納法により $L(H_n,(E_i,1))=\{w^*_i\}$ かつ $L(H_n,(E_i,\zeta))=L(G_n,E_i)\setminus\{w^*_i\}$ なので，$L(H_n)=M_n$．
$i<n$ については $\Lce^{n-i}x\Rce^{n-i}$ は $\Lce$ を含むので，それが $M_n$ に属するのは $L_n$ に属するとき，
すなわち $x\in L(G_n,E_i)$ のとき，かつそのときに限る．$F_1$ と $F_\zeta$ と交わると二つの成分が得られる．
$i=n$ については $\{(\eps,\eps)\}\Lc\cap F_\zeta=M_n$．
左辺が型 $1$ の生成規則は $n$ 個である．$i\ge2$ の各 $(E_i,\zeta)$ は二項の生成規則を 3 個，一項の生成規則を 2 個持ち，
$(E_1,\zeta)$ は $\Lce(E_0,1)\Rce$ だけを持つ（$(E_0,\zeta)$ は存在せず，対 $(1,1)$ は除かれる）．よって生成規則は $n+5(n-1)+1=6n-4$ 個である．

(c) 閉包は $M_n$ に相対的である．$i<n$ について $c_i=(\Lce^{n-i},\Rce^{n-i})$，$c_n=(\eps,\eps)$ と書く．
(b) により，$i<n$ について $\{c_i\}\Lc=L(G_n,E_i)$，$\{c_n\}\Lc=M_n=L(G_n,E_n)\setminus\{w^*_n\}$．
$(p_i,q_i)$ を $w^*_n$ の $E_i$-節点の文脈とする．したがって $p_iw^*_iq_i=w^*_n$ であり，$p_i,q_i$ は標識を持たない．

$0\le i<n$ について $\{w^*_i\}\RL=L(G_n,E_i)$ を示す．文脈 $c_i$ は $w^*_i$ を受け入れるので，$\{w^*_i\}\RL\subseteq\{c_i\}\Lc=L(G_n,E_i)$．
逆に $(p,q)\in\D(w^*_i)$，$y\in L(G_n,E_i)$ とする．補題~\ref{lem:nodes} により $pyq\in L_n$．
$pw^*_iq\in M_n$ は標識を持ち $w^*_i$ は持たないので，$p$ か $q$ が標識を持ち，$pyq\ne w^*_n$，すなわち $pyq\in M_n$．

$1\le i\le n$ について $\{u_i\}\RL=L(G_n,E_i)\setminus\{w^*_i\}$ を示す．
$(p,q)\in\D(u_i)$ で $y\in L(G_n,E_i)$ が標識を持つなら，補題~\ref{lem:nodes} により $pyq\in L_n$ であり，$pyq$ は標識を持つので $pyq\in M_n$．これで $\supseteq$ が得られる．
逆に，$\Lce^{n-i}u_i\Rce^{n-i}=u_n$ と $p_iu_iq_i$ は $L_n$ の標識付きの語なので，$c_i$ と $(p_i,q_i)$ は $u_i$ を受け入れる．
よって $\{u_i\}\RL\subseteq\{c_i\}\Lc\subseteq L(G_n,E_i)$ であり，$p_iw^*_iq_i=w^*_n\notin M_n$ なので，$\{u_i\}\RL$ のどの $x$ も $w^*_i$ ではない．

$w^*_i$ は $L(G_n,E_i)$ の型 $1$ のただ一つの語なので，$F_1$ と $F_\zeta$ と交わると，(b) の証明により
$\cl_1(\{w^*_i\})=\{w^*_i\}=L(H_n,(E_i,1))$，$\cl_\zeta(\{u_i\})=L(G_n,E_i)\setminus\{w^*_i\}=L(H_n,(E_i,\zeta))$．
核は型について均質なので，補題~\ref{lem:reduction} はこの型付き SPP を，同じ文法の型なしの SPP $(\{w^*_i\}\RL)$，$(\{u_i\}\RL)$ に変える．
\end{proof}

\begin{proposition}[観測器なしの枠]\label{prop:untyped-frames}
$n\ge1$ とする．$0\le d\le n-1$ と $\pi=\pi_1\cdots\pi_d\in\{l,r\}^d$ について，
$\Gamma\cup\{E_0,\dots,E_{n-1}\}$ 上の文字列 $T_d,\dots,T_0$ を，$T_d=\Lce E_{n-d-1}\Rce$，および $1\le j\le d$ について
\[
T_{j-1}=\Lfl T_jE_{n-j}\Rfl\ \text{（}\pi_j=l\text{ のとき）},\qquad
T_{j-1}=\Lfl E_{n-j}T_j\Rfl\ \text{（}\pi_j=r\text{ のとき）}
\]
で定め，$T_\pi=T_0$ と書く．
$F_n$ を，非終端記号 $E_0,\dots,E_{n-1}$ と $G_n$ におけるそれらの生成規則，新しい初期記号 $S$，
およびそのようなすべての $\pi$ についての生成規則 $S\to T_\pi$ からなる文法とする．
このとき $L(F_n)=M_n$ であり，$M_n$ に相対的に，$F_n$ は $E_k$ に文脈 $(\Lce^{n-k},\Rce^{n-k})$，$S$ に文脈 $(\eps,\eps)$ を用いて強 $1$-FCP を持つ．
非終端記号は $n+1$ 個，生成規則は $2^n+2n-2$ 個，$\rho=3n$ である．
\end{proposition}
\begin{proof}
非終端記号 $E_k$ は $L(G_n,E_k)$ を生成する．$d-j$ についての帰納法により，$T_j$ から導出されるどの終端語も
$L(G_n,E_{n-j})$ に属し $\Lce$ を含む．よって $L(F_n)\subseteq M_n$．
逆に $m\in M_n$ とする．$G_n$ における $m$ の導出木は生成規則 $E_i\to\Lce E_{i-1}\Rce$ を用いる節点を持つ．
そのような節点のうち深さ $d$ が最小のものを $\nu$ とする．
このとき $d\le n-1$ であり，$\nu$ の真の祖先はすべて二項の生成規則を用いる．
$\nu$ の祖先で $\nu$ への道が左右どちらに進むかを記録すると $\pi\in\{l,r\}^d$ が得られ，
$m$ は $T_\pi$ から導出される．道から外れた部分木と $\nu$ の下の部分木は $E_k$ から導出される．
$k<n$ について，命題~\ref{prop:deletion-upper}(b) の証明により，$M_n$ に相対的に
$\{(\Lce^{n-k},\Rce^{n-k})\}\Lc=L(G_n,E_k)=L(F_n,E_k)$ であり，$\{(\eps,\eps)\}\Lc=M_n=L(F_n,S)$．
$S$ の生成規則は $\sum_{d<n}2^d=2^n-1$ 個，$E_0$ の生成規則は 1 個，$1\le k<n$ の各 $E_k$ の生成規則は 2 個であり，
$|T_\pi|=3d+3\le3n$．
\end{proof}

\begin{lemma}[核の捕獲]\label{lem:capture}
$L'$ を言語，$z=RT\in L'$ とし，語 $T^\circ$ が
\begin{equation}\label{eq:capture}
\{R':\ R'T\in L'\ \text{かつ}\ R'T^\circ\notin L'\}=\{R\}
\end{equation}
を満たすとする．$G$ を $L(G)=L'$ で，$L'$ に相対的な SPP $(K_A\RL)_A$ を持つ CFG とする．
$z$ のある導出木で，ラベル $B$ の節点が $R$ の中にある収量の出現 $x$ を持つならば，$x\in K_B$ である．
\end{lemma}
\begin{proof}
$z=pxq$ と書く．このとき $(p,q)\in\mathcal C(G,B)$ であり，補題~\ref{lem:reach} により $pK_Bq\subseteq pK_B\RL q\subseteq L'$．
$x$ は $R$ の中にあるので，ある $q_0$ について $q=q_0T$ である．$q'=q_0T^\circ$ とおくと，\eqref{eq:capture} により $pxq'=RT^\circ\notin L'$．
$x\in L(G,B)\subseteq K_B\RL$ なので $K_B\R\subseteq\D(x)\not\ni(p,q')$ であり，ある $e\in K_B$ について $peq'\notin L'$．
$peq=(peq_0)T\in L'$ なので，\eqref{eq:capture} から $peq_0=R$，すなわち $e=x$．
\end{proof}

\begin{lemma}[一語の削除]\label{lem:one-word}
$L$ を言語，$w=RT^\circ\in L$ とし，$T\ne T^\circ$ を，$RT\in L$ であり，すべての語 $y$ について $yT\in L$ ならば $yT^\circ\in L$ となる語とする．
このとき $z=RT$ と $T^\circ$ は $L'=L\setminus\{w\}$ について \eqref{eq:capture} を満たす．
$n\ge2$ のとき，これは $L=L_n$，$w=w^*_n$，$T^\circ=\Lfl\ta\ta\Rfl\Rfl^{n-1}$，$T=\Lce\ta\Rce\Rfl^{n-1}$ に適用でき，
$|R|=|w^*_n|-(n+3)=3\cdot2^n-n-5$ である．
\end{lemma}
\begin{proof}
$T\ne T^\circ$ なので $z\ne w$，よって $z\in L'$．$R'T\in L'$ ならば $R'T^\circ\in L$ であり，
$R'T^\circ\notin L'$ となるのは $R'T^\circ=w$，すなわち $R'=R$ のとき，かつそのときに限る．
$L_n$ について：$w^*_n$ の最後の $E_1$-節点は $\Lfl\ta\ta\Rfl$ であり，その $n-1$ 個の祖先の閉じ括弧が続くので，$w^*_n$ は $T^\circ$ で終わる．
$RT\in L_n$ はこの節点を $E_1$ の語 $\Lce\ta\Rce$ に置き換えたものである．
$yT\in L_n$ ならば，$yT$ の中の $\Lce\ta\Rce$ の出現は $E_1$-節点であり（補題~\ref{lem:nodes}），それを $\Lfl\ta\ta\Rfl$ に置き換えると $yT^\circ\in L_n$．
\end{proof}

\begin{theorem}[核の下界]\label{thm:general-lower}
$L'$ を有限言語とし，$z=RT\in L'$ と $T^\circ$ が \eqref{eq:capture} を満たすとする．
$G=(N,\Sigma,P,I)$ を，$L(G)=L'$ で，有限の核の SPP $(K_A\RL)_A$ を持つ任意の CFG とし，
$\rho$ を右辺の記号の最大数，$\kappa=\max\{|e|:e\in\bigcup_AK_A\}$ とする．このとき $\rho\,|N|\,\max(\kappa,1)\ge|R|$．
\end{theorem}
\begin{proof}
$R=\eps$ なら示すことはない．$T=\eps$ かつ $R\ne\eps$ なら，$z$ の導出木の根の収量 $z=R$ は $R$ の中にあるので，
補題~\ref{lem:capture} によりその根のラベル $S$ について $z\in K_S$ であり，$\kappa\ge|R|$．
以下 $R,T\ne\eps$ とし，$z$ の導出木 $\mathcal T$ を，節点の数が最小のものにとる．$\mathcal T$ の道の上で非終端記号は繰り返さない．
$A$ がある節点とその真の子孫のラベルで，収量がそれぞれ $\alpha y\beta$ と $y$ ならば，
$\alpha\beta\ne\eps$ からは無限個の語 $u\alpha^ty\beta^tv\in L(G)$ が得られて $L'$ の有限性に反し，
$\alpha\beta=\eps$ ならより小さな木がとれるからである．
収量が $z$ の位置 $|R|$ と $|R|+1$ の両方を含む節点を\emph{鎖節点}と呼ぶ．鎖節点は根からの道をなすので，高々 $|N|$ 個である．
各位置 $j\le|R|$ について，$j$ を含む最も深い鎖節点は，$j$ を右辺の終端記号として持つか，$j$ を含む鎖節点でない子を持つ．
そのような子は $R$ の中にある．
よって $R$ は高々 $\rho|N|$ 個の断片で覆われ，各断片は一つの終端記号か $R$ の中にある節点の収量である．
補題~\ref{lem:capture} により，そのような収量は核に属するので，長さは $\kappa$ 以下である．
したがって $|R|\le\rho|N|\max(\kappa,1)$．
\end{proof}

\begin{theorem}[$M_n$ の核の下界]\label{thm:deletion}
$n\ge2$ とする．有限の核の SPP $(K_A\RL)_A$ を持つ，またはある型付けについて型付き極弱 FKP を持つ $M_n$ のどの CFG も，
\[
\rho\,|N|\,\max(\kappa,1)\ \ge\ 3\cdot2^n-n-5
\]
を満たす．
\end{theorem}
\begin{proof}
補題~\ref{lem:one-word} と定理~\ref{thm:general-lower} による．
\end{proof}

下界はパラメータのあいだの交換条件であり，二つの逆の仕方で $O(n)$ 倍の範囲で最適である．
命題~\ref{prop:deletion-upper}(c) の提示は $\rho$ と $|N|$ が小さく，$\rho|N|\kappa\le8n\max(3\cdot2^{n-1}-2,2n+1)$ を満たす．
非終端記号 $S$ 一つと，各 $w\in M_n$ についての生成規則 $S\to w$ と，核 $K_S=\{u_n\}$ からなる列挙の文法は，
強 $1$-FKP を持ち（$u_n$ の文脈は $(\eps,\eps)$ だけなので $\{u_n\}\RL=M_n$），$|N|$ と $\kappa$ が小さく，
$\rho|N|\kappa=(3\cdot2^n-3)(2n+1)$ だが，二重指数個の生成規則を持つ．
したがってどのパラメータもすべての提示で指数的になるわけではない．定理が排除するのは，すべてのパラメータが小さい提示である：
$\rho|N|\le p(n)$ ならば $\kappa\ge(3\cdot2^n-n-5)/p(n)$．


\begin{remark}[標識の観測器と $L_n$ の見本]\label{rem:Ln-sU}
型付きの核は型について均質なので，観測器は $L_n$ の見本の側に不利に働きうる．
$L$ を有限とし，$(G,\tau)$ を型付き SPP $(\cl_{\tau(A)}(K_A))_A$ を持つ $L$ の $s$-型付き CFG とする．
$\mu\in\widehat s(L)$ について，型 $\mu$ のある初期記号 $S$ は $\cl_\mu(K_S)\ne\emptyset$ を満たすので（有限な $L$ では $\emptyset\RL=\emptyset$ なので）$K_S\ne\emptyset$ であり，健全性により $K_S\subseteq L\cap F_\mu$．
よって $\kappa\ge\max_{\mu\in\widehat s(L)}\min\{|w|:w\in L\cap F_\mu\}$ であり，これは
\citet[式~($\spadesuit$)，Cor.~6.6]{kuriyama-tcs1} が型付き文法の収量を下から抑えるのに用いた量である．
$s_U$ のもとの $L_n$ では $L_n\cap F_1=\{w^*_n\}$ なので，$L_n$ のどの $s_U$-型付きの見本の提示も $w^*_n$ を挙げ，$\kappa\ge3\cdot2^n-2$ である．
一方，$G_n$ を補題~\ref{lem:product} と命題~\ref{prop:duality}(c) で細分化すると，
文脈の長さが $2n$ 以下の，大きさ $O(n)$ の $s_U$-型付き強 $1$-FCP の文法が得られる．
この不利は観測器による：$L(G_n,E_i)$ のすべての語は $L_n$ で同じ文脈を持つので（補題~\ref{lem:nodes}），$G_n$ を長さの $2$ を法とする剰余で細分化すると，核 $u_i$ と $\Lfl u_{i-1}u_{i-1}\Rfl$ を持つ大きさ $O(n)$ の強 $1$-FKP の文法が得られる．
\end{remark}

%======================================================================
\section{学習アルゴリズムへの帰結}\label{sec:applications}

この節では定理~\ref{thm:deletion} の下界を，著者の正例からの再構成と \citet{ky2017} の主学習器に対する下界に変える．
枠が短い双対学習器には，別の反駁の論法が要る（命題~\ref{prop:dual-refutation}）．
学習器は \citet{ky2017} の Algorithms~1 と~2 であり，Clark のもとの規則の検査とその修正については \citet{leiss2014} を参照．

\subsection*{正例からの学習}

\citet[Sec.~4.1]{kuriyama-tcs1} を思い出すと，$\mathsf{Rec}_s(K)$ は，有限の標本 $K$ の語のすべての非空因子 $x$ についての非終端記号 $[\![x]\!]$，
開始記号 $\widehat S$，および次の生成規則を持つ：
(B) $[\![xy]\!]\to[\![x]\!][\![y]\!]$，
(U) $x\approx_{K,s}y$（すなわち $s(x)=s(y)$ かつ $\D_K(x)\cap\D_K(y)\ne\emptyset$．ここで $\D_K$ は $K$ に相対的な $\D$）なる $[\![x]\!]\to[\![y]\!]$，
(L) $[\![a]\!]\to a$，
(S) $w\in K\cap\Sigma^+$ についての $\widehat S\to[\![w]\!]$，および $\eps\in K$ なら $\widehat S\to\eps$．

\begin{lemma}[再構成された文法は核の SPP を持つ]\label{lem:rec-kernel}
$s$ を型付け，$L\in\Ccf s$ を $\eps\notin L$ なる空でない言語とし，$K\subseteq L$ を $L(\mathsf{Rec}_s(K))=L$ となる有限集合とする．
$\|K\|=\sum_{w\in K}(|w|+1)$ と書く．
$G_K$ を，$\mathsf{Rec}_s(K)$ から $\widehat S$ とその生成規則を除き，$w\in K$ なるすべての $[\![w]\!]$ を初期記号と宣言して得られる文法とする．
このとき $L(G_K)=L$ であり，$L$ に相対的に，$(\{x\}\RL)_{[\![x]\!]}$ は $G_K$ の SPP である．
よって $G_K$ は，$\rho_{G_K}\le2$，$|N_{G_K}|\le\|K\|^2/2$，$\kappa_{G_K}\le\max_{w\in K}|w|$ で極弱 $1$-FKP を持つ．
\end{lemma}
\begin{proof}
これは \citet[Thm.~4.2]{kuriyama-tcs1} の健全性の不変式を閉包の形で述べたものである．$\widehat S$ はどの右辺にも現れず，$\eps\notin K$ なので，$L(G_K)=L(\mathsf{Rec}_s(K))=L$．
前不動点であること：(B) については補題~\ref{lem:clark} により $\{xy\}\RL\supseteq\{x\}\RL\{y\}\RL$．
(L) については $a\in\{a\}\RL$．
(U) については，$K\subseteq L$ から $\D_K(x)\subseteq\D_L(x)$ が得られるので，$x$ と $y$ は等しい型と $L$ での共通の文脈を持ち，
$\sim_s$-代入可能性から $\D_L(x)=\D_L(y)$，すなわち $\{x\}\RL=\{y\}\RL$．
健全性：$w\in K\subseteq L$ について $(\eps,\eps)\in\D_L(w)$ なので $\{w\}\RL\subseteq L$．
一方，初期非終端記号にわたる和は $L(G_K)=L$ を含む．
大きさの評価：どの右辺も高々 2 個の記号を持つ．語 $w$ は高々 $|w|(|w|+1)/2$ 個の非空因子を持つ．核は観測された一つの因子である．
\end{proof}

\begin{corollary}[一語を削除した標的の特性データ]
\label{cor:rec-data}
$n\ge2$ とし，$s$ を $M_n\in\Ccf s$ となる任意の型付けとする．
たとえば $s_U$ \citep[Prop.~6.4(iv)]{kuriyama-tcs1}，または $s_U$ の任意の細分化である．
$K\subseteq M_n$ が有限で $L(\mathsf{Rec}_s(K))=M_n$ ならば
\[
\|K\|^3\ \ge\ 3\cdot2^n-n-5 .
\]
特に，\citet[Sec.~2]{kuriyama-tcs1} の意味での $\mathsf{Rec}_s$ に対する $M_n$ のどの特性標本 $\mathcal F$ も
$\|\mathcal F\|\ge(3\cdot2^n-n-5)^{1/3}$ を満たす．
\end{corollary}
\begin{proof}
補題~\ref{lem:rec-kernel} の $G_K$ に定理~\ref{thm:deletion} を適用する：
$3\cdot2^n-n-5\le\rho|N|\max(\kappa,1)\le2\cdot(\|K\|^2/2)\cdot\|K\|$．
特性標本 $\mathcal F$ は $L(\mathsf{Rec}_s(\mathcal F))=M_n$ を満たす（$K=\mathcal F$ とおく）．
$s_U$ の細分化 $h$ については，$h$-型が等しければ $s_U$-型も等しいので，$\sim_{s_U}$-代入可能性から $\sim_h$-代入可能性が従い，$M_n\in\Ccf h$．
\end{proof}

\begin{remark}[固定型付けの論文との関係]\label{rem:rec-data}
\citet[Prop.~6.4]{kuriyama-tcs1} は，$\mathsf{Rec}_{s_U}$ に対する $L_n$ のどの特性標本も $w^*_n$ を含むことを示すが，$M_n$ の評価は与えていない．
系~\ref{cor:rec-data} は，$M_n\in\Ccf s$ となるすべての $s$ についてそれを与える（$n\ge2$ について $M_n\notin\Ccf{s_0}$ なので $s_0$ は除かれる）．指数 $1/3$ が最適であるとは主張しない．
$M_n$ のどの特性標本も長さ $2^{\Omega(n)}$ の語を含むかどうかは未解決である．
\end{remark}

\subsection*{所属質問を用いる学習器}

$M_n$ は有限なので，どちらの学習器も $M_n$ のどの正の提示の上でも収束し（すべての語が現れた後は集合は高々もう一度しか変わらない），
命題~\ref{prop:learners} により極限の文法は正しい．問題は，どれだけのデータが要るかである．

\begin{corollary}[主学習器の収束時のデータ]\label{cor:primal-data}
$\mathrm{Primal}_{r,k}$（任意の $s$ について型付き，または型なし）が $M_n$ の正の提示の上で段階 $i$ で収束するならば，
それまでに見たデータ $D_i$ は
\[
\|D_i\|\ \ge\ \Bigl(\frac{3\cdot2^n-n-5}{2r+1}\Bigr)^{1/(2k+2)},
\qquad \|D\|=\textstyle\sum_{d\in D}(|d|+1)
\]
を満たす．
\end{corollary}
\begin{proof}
極限の文法 $G=G_i$ は，その非終端記号で添字づけられた SPP $(K\RL)_K$ を持つ（事実~\ref{fact:ky}(b)，命題~\ref{prop:learners}）．
その核は $D_i$ の高々 $k$ 個の部分文字列の集合なので，$\kappa\le\|D_i\|$ かつ $|N|\le|\Sub(D_i)|^k\le\|D_i\|^{2k}$．
その右辺は高々 $r$ 個の非終端記号とデータの語の $r+1$ 個の因子からなるので，$\rho\le(r+1)\|D_i\|+r\le(2r+1)\|D_i\|$．
定理~\ref{thm:deletion} を適用する．
\end{proof}
$r,k$ を固定すれば系~\ref{cor:primal-data} の下界は $2^{\Omega(n)}$ である．$k$ が $n$ とともに増える場合は，示した下界だけを主張する．

\begin{proposition}[双対学習器の反駁の下界]
\label{prop:dual-refutation}
$n\ge2$，$r\ge2$，$k\ge1$ とする．型なしの双対学習器 $\mathrm{Dual}_{r,k}$，または型付きの双対学習器 $\mathrm{Dual}^{s_U}_{r,k}$ が，
$M_n$ の正の提示の上で段階 $i$ で収束するならば，$w^*_{n-1}$ は $D_i$ のある語の因子である．
特に，あるデータは長さ $3\cdot2^{n-1}-2$ 以上を持つ．
\end{proposition}
\begin{proof}
$t\le i$ を $J$ と $H$ が最後にリセットされた段階とする．したがって極限の文法は
$J_t=\Con(D_t)$，$H_t=\Sub^{\le r+1}(D_t)$ から作られ，$\mathcal E_i=\Sub(D_i)$ の上で検査される．
\citet[Thm.~3]{ky2017} の証明（および命題~\ref{prop:learners}）により，極限の文法のどの生成規則も妥当なので，
妥当でない候補の生成規則はどれも $\mathcal E_i$ の上の検査で落ちなければならない．

$M_n$ のどの語も $\Lce y\Rce$ か，$y,y'$ の少なくとも一方が型 $\zeta$ である $\Lfl yy'\Rfl$ である．
$D_t$ のあるデータについて，$c=(\Lce,\Rce)$，$y'$ が型 $\zeta$ の $c=(\Lfl,y'\Rfl)$，または $y$ が型 $\zeta$ の $c=(\Lfl y,\Rfl)$ をとる．
どの場合も $c\in J_t$ であり，$M_n$ に相対的に $\{c\}\Lc=L(G_n,E_{n-1})$．
極限の文法は $\Lfl$ を含む語を生成するので，$D_t$ のあるデータは $\Lfl$ を含み，したがってそれより後の $\Rfl$ も含むので，
$(\Lfl,\eps,\Rfl)\in H_t$．

型なしの場合：候補 $\{(\eps,\eps)\}\to\Lfl\{c\}\{c\}\Rfl$ は，$w^*_n=\Lfl w^*_{n-1}w^*_{n-1}\Rfl\notin M_n=\{(\eps,\eps)\}\Lc$ なので妥当でない．
これが $\mathcal E_i$ の上の検査で落ちるのは，$\Lfl x_1x_2\Rfl\notin M_n$ となる $x_1,x_2\in \mathcal E_i\cap\{c\}\Lc$，すなわち $x_1=x_2=w^*_{n-1}$ があるときに限る．

型付きの場合：型の整合する候補 $(\{c\},\zeta)\to\Lfl(\{c\},1)(\{c\},\zeta)\Rfl$ は妥当でない．
$\Lfl w^*_{n-1}y\Rfl$（$y\in L(G_n,E_{n-1})$ は型 $\zeta$）は入れ子の深さが $n$ なので $\{c\}\Lc$ に属さないからである．
検査が落ちうるのは $x_1\in \mathcal E_i\cap\{c\}\Lc\cap F_1=\mathcal E_i\cap\{w^*_{n-1}\}$ のときだけである．
\end{proof}

これは，候補の生成規則を追加で刈り込まない \citet{ky2017} の Algorithm~1 についてのものである．$H_n$ の全長は $O(n^2)$ である（命題~\ref{prop:mark-total}）にもかかわらず，$\mathrm{Dual}^{s_U}_{r,k}$ はこの長いデータを要する．同じ論法は \citet{kanazawa2023wollic} の Algorithm~3（演算に共通部分を含み，$r\ge2$，$s\ge1$）にそのまま通る：収束時には商が $L(G_n,E_{n-1})$ の原子 $e$ が非終端記号であり，初期の原子から $\Lfl e\,e\Rfl$ への生成規則は $w^*_{n-1}$ によってしか棄却されない．拡張正則の演算のもとでは，$H_n$ は商 $\{c_i\}\Lc$ と $\{c_i\}\Lc\cap\Gamma^*\Lce\Gamma^*$（$c_i$ は命題~\ref{prop:deletion-upper}(c) の証明のもの）による SPP を持ち，Kanazawa の正規化のあとでも，その原子はどれも長さ $(n+1)^2$ 以下の一つの語に現れ，$k\ge8$ なら Algorithm~3 は $M_n$ を同定する \citep[Thm.~8]{kanazawa2023wollic}．したがって，\citet{kanazawa2023wollic} のデータ効率の性質 (d$^{\dagger\prime}$) は追加の例の個数を抑えるが，その長さまで目標の文法の大きさの多項式で抑えることはできない．

\begin{remark}[固定されたすべての観測器に対する一つの主学習器]\label{rem:uniform-known}
$\bigcup_s\Ccf s$ のどの言語も，極弱 $1$-FKP を持つ階数 $2$ の CFG を持つ（\citealp[Prop.~9.7]{kuriyama-tcs1} と \citealp[Sec.~6.3]{yoshinaka2015}．直接には，$L\setminus\{\eps\}$（空でなければ）の特性標本に補題~\ref{lem:rec-kernel} を適用し，$\eps\in L$ なら核 $\{\eps\}$ を持つ初期の $A\to\eps$ を加える）．
よって型なしの学習器 $\mathrm{Primal}_{2,1}$ は，所属質問を用いてこの和のすべての言語を同定する（事実~\ref{fact:ky}(b)，\citealp[Thm.~6]{ky2017}）．
この和は正例だけからは同定できない（第~\ref{sec:intro} 節）．
\end{remark}


%======================================================================
\section{観測器の役割}
\label{sec:duality}

次の観察は，\citet[Sec.~4]{ky2017} の型なしの規則の妥当性の同値と，制限したガロア接続から従う．
\begin{proposition}[文脈と核の双対性]\label{prop:duality}
$s\preceq h$ を $\Sigma^+$ の型付けとし，$k\ge1$ とする．
\begin{enumerate}[label=\textup{(\alph*)},nosep]
\item $\tFKP{k}{vw}{s}=\tFKP{k}{w}{s}$．
\item $\tFKP{k}{vw}{h}\subseteq\tFKP{k}{vw}{s}\subseteq\FKP{k}{vw}=\FKP{k}{w}$．
\item $\FCP{k}{vw}\subseteq\tFCP{k}{vw}{s}\subseteq\tFCP{k}{vw}{h}$ かつ
$\FCP{k}{st}\subseteq\tFCP{k}{st}{s}\subseteq\tFCP{k}{st}{h}$．
\end{enumerate}
\end{proposition}
\begin{proof}
(a) 包含 $\supseteq$ について，$(\cl_{\tau(A)}(L(G,A)))_A$ はどの $s$-型付き $(G,\tau)$ の SPP でもあることに注意する：
生成規則 $A\to w_0A_1\cdots A_nw_n$ について，集合 $P=w_0\cl(L(G,A_1))w_1\cdots\cl(L(G,A_n))w_n$ は $F_{\tau(A)}$ に含まれ，
補題~\ref{lem:star} により $P\R=(w_0L(G,A_1)w_1\cdots L(G,A_n)w_n)\R\supseteq L(G,A)\R$ なので，
補題~\ref{lem:typed-galois} により $P\subseteq\cl_{\tau(A)}(L(G,A))$．
$S\in I$ については $(\eps,\eps)\in L(G,S)\R$ から $\cl(L(G,S))\subseteq L$．
$\subseteq$ について，$(G,\tau)$ が型付き SPP $(\cl_{\tau(A)}(K_A))_A$ を持つとする．
補題~\ref{lem:reduction} により $(K_A\RL)_A$ は型なしの SPP である．
$w\in K_A$ について $A\to w$ を加えて得られる事実~\ref{fact:ky}(a) の文法 $G'$ は，
$K_A\subseteq F_{\tau(A)}$ なので依然として $s$-型付きであり，$L(G')=L$ かつ $L(G',A)\RL=K_A\RL$ を満たす．
$F_{\tau(A)}$ と交わると $\cl_{\tau(A)}(L(G',A))=\cl_{\tau(A)}(K_A)$ が得られる．

(b) $(G,\tau)$ を型付き SPP $(\cl^h_{\tau(A)}(K_A))_A$ を持つ $h$-型付き文法とする．
（$h$ についての）補題~\ref{lem:reduction} により $(K_A\RL)_A$ は型なしの SPP である．これは $\FKP{k}{vw}$ への所属を示す．
さらに $(G,\pi\circ\tau)$ は $s$-型付きで，$K_A\subseteq F^h_{\tau(A)}\subseteq F^s_{\pi\tau(A)}$ なので，
（$s$ についての）補題~\ref{lem:reduction} は同じ型なしの SPP を型付き SPP $(\cl^s_{\pi\tau(A)}(K_A))_A$ に変える．
等号は事実~\ref{fact:ky}(a) である．

(c) $G$ を，SPP $(X_A)_A$，$X_A=C_A\Lc$ または $X_A=C_A\Lc\cap F^s_{\sigma(A)}$ を持つ型なしまたは $s$-型付きの文法とする．
補題~\ref{lem:product} のように $G^h$ を作り，$X_{(A,\nu)}=X_A\cap F^h_\nu=C_A\Lc\cap F^h_\nu$ とおく
（型付きの場合は $F^h_\nu\subseteq F^s_{\sigma(A)}$）．
$G^h$ の生成規則について，右辺 $w_0X_{(A_1,\nu_1)}\cdots w_n$ は $X_A$ に含まれ（$G$ の前不動点），$F^h_\nu$ にも含まれる（型は掛け合わされる）ので，
この族は前不動点である．$\bigcup_\nu X_S\cap F^h_\nu=X_S$ なので健全である．
$X_A=L(G,A)$ として同じ計算を行い補題~\ref{lem:product} を用いると，強い場合が得られる．
\end{proof}

したがって観測器は主学習器のクラスを決して広げない：$\mathrm{Primal}^s_{r,k}$ が $L_*$ のある正の提示の上で $L_*$ の文法に収束するなら，その文法は階数 $r$ 以下で極弱 $k$-FKP を持つ（命題~\ref{prop:learners}）ので，$\mathrm{Primal}_{r,k}$ は $L_*$ のすべての正の提示の上で収束する（事実~\ref{fact:ky}(b)）．

\subsection*{2 元の半束での真包含}

命題~\ref{prop:duality} の型付きの包含は，すべての $k\ge1$ について，すでに $s_0$ と $U_1$ への指示型付けのあいだで真である．
以下の証人は，一方の側で強 $1$-性質を持ち，他方の側でどの $k$ についても極弱の性質を持たない．
型なしのクラスと $s_0$-型付きのクラスのあいだの端点の包含については，真であるとは主張しない．

\begin{proposition}[観測は文脈を助ける]\label{prop:strict-context}
$\Sigma=\{a,b,c,d\}$ とし，
$L=\{a^mcb^{2m}:m\ge0\}\cup\{a^mdb^n:n\le m\le2n\}$
（\citealp[Prop.~2(i)]{ky2017}）とする．このとき $L\in\tFCP1{st}{h_{\{c\}}}$ だが，
$L\notin\FCP{}{vw}$ かつ $L\notin\tFCP{}{vw}{s_0}$．
\end{proposition}
\begin{proof}
$h=h_{\{c\}}$ とする．非終端記号 $S_1$（型 $\zeta$）と $S_2$（型 $1$）をともに初期記号とし，
生成規則 $S_1\to aS_1bb\mid c$，$S_2\to aS_2b\mid aaS_2b\mid d$ を持つ文法は $h$-型付きであり，
$L(G,S_1)=L\cap F_\zeta$，$L(G,S_2)=L\cap F_1$．$\{(\eps,\eps)\}\Lc=L$ なので，
どちらの成分も $\{(\eps,\eps)\}\Lc\cap F_\mu$ であり，$G$ は型付き強 $1$-FCP を持つ．

$L\notin\FCP{}{vw}$ は \citet[Prop.~2(i)]{ky2017} であり，その議論は $s_0$-型付きの文法も除外する．
$L$ の任意の CFG（型付きでもそうでなくても）について，$a^pcb^{2p}$ をポンプすると，初期記号 $S$ と非終端記号 $E$ で
$S\Rightarrow^*a^{i_1}Eb^{j_1}$，$i>0$ なる $E\Rightarrow^+a^iEb^{2i}$，$E\Rightarrow^*a^{i_2}cb^{j_2}$ となるものが得られる．
議論が用いるのは，SPP の $E$-成分が $C_E$ は有限，$F\in\{\Sigma^*,\Sigma^+\}$ として $C_E\Lc\cap F$ であることだけである：
その文脈は $(a^{n-i_2},b^{2n-j_2})$ の形なので，$m$ が大きければ空でない語 $a^{3m}db^{2m}$ がそれに属し，
到達文脈 $(a^{i_1+ti},b^{j_1+2ti})$ に補題~\ref{lem:reach} を適用するとすべての $t$ について $a^{i_1+ti+3m}db^{j_1+2ti+2m}\in L$ となるが，これは $t$ が大きいと成り立たない．
\end{proof}

\begin{proposition}[観測は核に不利に働く]\label{prop:strict-kernel}
$\Sigma=\{a,b,y,z\}$，$P=\{a^mb^n:m\ne n\}$ とし，
\[
L_K=P\cup\{a^mzb^n:n\le m\}\cup\{a^myb^n:m\le n\}
\]
とする．このとき $L_K\in\tFKP1{st}{s_0}\subseteq\FKP1{vw}$ だが，すべての $k$ について $L_K\notin\tFKP{k}{vw}{h_{\{y,z\}}}$．
\end{proposition}
\begin{proof}
\emph{上界.}  $s_0$-型がともに $0$ の二つの初期非終端記号 $S_+,S_-$ と，生成規則
\[
S_+\to a\mid z\mid aS_+\mid aS_+b,\qquad
S_-\to b\mid y\mid S_-b\mid aS_-b
\]
を用いる．それぞれの言語は
\[
L_+=\{a^mb^n:m>n\}\cup\{a^mzb^n:m\ge n\},\qquad
L_-=\{a^mb^n:m<n\}\cup\{a^myb^n:m\le n\}
\]
であり，その和は $L_K$ である．$L_K$ に相対的に，
$\D(z)=\{(a^p,b^q):p\ge q\}$，$\D(y)=\{(a^p,b^q):p\le q\}$．
$\{z\}\RL=L_+$ を示す．$\{y\}\RL=L_-$ の証明は対称である．
語 $x$ が $\{z\}\RL$ に属するのは，すべての $p\ge q\ge0$ について $a^pxb^q\in L_K$ となるとき，かつそのときに限る．$p=q=0$ とすると $x\in L_K$．
$x=a^mb^n$（$m\ne n$）なら，すべての $p\ge q$ について $a^{m+p}b^{n+q}\in L_K$ となるのは，すべての $p-q\ge0$ について $m+p\ne n+q$ となるとき，すなわち $m>n$ のときに限る．
$x=a^mzb^n$（$n\le m$）なら，すべての $p\ge q$ について $n+q\le m+p$．
$x=a^myb^n$ なら，$p>n-m$ について $a^{m+p}yb^{n}\notin L_K$．
よって $\{z\}\RL=L_+$．したがって生成される二つの成分はどちらも一元の核の閉包である．
どちらも $\eps$ を含まないので，この文法は $s_0$-型付き強 $1$-FKP を持つ．

\emph{下界.}  $h=h_{\{y,z\}}$ とし，$L_K$ の $h$-型付き $(G,\tau)$ が型付き SPP $(\cl_{\tau(A)}(K_A))_A$ を持つと仮定する．
$P$ のどの語も型 $1$ を持つので，$P\subseteq\bigcup\{\cl_1(K_S):S\in I,\ \tau(S)=1\}$．
そのような $S$ について，健全性により $K_S\subseteq F_1\cap\cl_1(K_S)\subseteq F_1\cap L_K=P$．
$p$ を，有限個のそのような $K_S$ に属するすべての $a^mb^n$ にわたる $|m-n|$ の上界とする
（それらがすべて空なら $p=0$．ここでは空の核の閉包は空である．$(ba,\eps)$ はどの語の文脈でもないからである）．
$|i-j|>p$ ならば $(a^i,b^j)\in K_S\R$ である．$|m-n|\le p$ について $a^{i+m}b^{n+j}\in P$ だからである．
$|m-n|>p$ なる語 $a^mb^n$ は $i-j=n-m$ なる文脈 $(a^i,b^j)$ を満たさないので，$K_S\RL$ に属さない．
よって和は $P$ の語 $a^{p+1}$ を取りこぼし，矛盾である．
この相殺の議論は \citet[Prop.~2(ii)]{ky2017} のものである．
\end{proof}

命題~\ref{prop:duality}(b) により，核の失敗はどの観測器のもとでも残る：強 $1$-FCP を持つが極弱 FKP を持つ CFG を持たない言語
$\{a^mb^n:m\ne n\}$ \citep[Prop.~2(ii)]{ky2017} は，どの型付け $s$ についても，型付き極弱 FKP を持つ $s$-型付き CFG を持たない．
弱 FCP は，命題~\ref{prop:duality}(c) の強 FCP とも極弱 FCP とも異なる振る舞いをする：細分化について単調でない．
\citet[Secs.~5--6]{ky2017} の言語を用いる．$\Sigma_0=\{a,b,c,d,e,f,\#,\$,\%\}$ の上で，$V=\{a,b,\#\}^*$，$V_0=\{a,b\}^*$ として
\begin{align*}
L_1&=\{v\$z: v\in V,\ z=(\mathrm{erase}_\#(v))^R\},\\
L_2&=\{wyc^id^ie^jz: w,z\in V_0,\ y\in(\#V_0)^*,\ i,j\ge0,\ |w|_a\ge|w|_b\},\\
L_3&=\{wyc^id^je^jz: w,z\in V_0,\ y\in(\#V_0)^*,\ i,j\ge0,\ |w|_a\le|w|_b\},\\
L_4&=\{v\$zc^k\%f^l: v\in V,\ z\in V_0,\ k,l\ge0\},\\
L_5&=\{vc^id^ie^jzc^k\%f^l: v\in V,\ z\in V_0,\ i,j,k,l\ge0,\ j\ne k\}
\end{align*}
とし，$L_{123}=L_1\cup L_2\cup L_3$，$L'=L_{123}\cup L_4\cup L_5$ とする．
また $J_0=\{vc^id^ie^iz: v\in V,\ z\in V_0,\ i\ge0\}$ とおく．
$V\subseteq L_2\cup L_3$ かつ $J_0\subseteq L_2\cup L_3$ に注意する．

\begin{proposition}[弱 FCP は単調でない]\label{prop:weak-not-preserved}
$L_1^@$ を，$L_1$ の $\$$ を新しい文字 $@$ に置き換えたものとし，
$\Sigma=\Sigma_0\cup\{@\}$，$L^+=L'\cup L_1^@$，$h=h_{\{@\}}$ とおく．したがって $s_0\preceq h$ である．このとき：
\begin{enumerate}[label=\textup{(\alph*)},nosep]
\item $L^+\in\FCP2{w}$ かつ $L^+\in\tFCP2{w}{s_0}$；
\item すべての $k$ について $L^+\notin\tFCP{k}{w}{h}$．
\end{enumerate}
よって，すべての $k\ge2$ について $\tFCP{k}{w}{s_0}\not\subseteq\tFCP{k}{w}{h}$．
\end{proposition}
\begin{proof}
証明は \citet[Lemmas~8 and~12]{ky2017} の文脈集合を新しい文字 $@$ に照らして確かめるものであり，付録~\ref{app:weak} で与える．
\end{proof}

このように，強 FCP や極弱 FCP と違って，型付き弱 FCP は細分化について単調でなく，それはすでに $s_0$ と 2 元の指示型付けのあいだで起こる．
$k=1$ の場合は未解決として残る（未解決問題~\ref{op:weak-fcp}）．

これに対し，命題~\ref{prop:strict-context} の文脈の側の失敗は，1 ビットの観測で修復される．

\subsection*{枠を助ける観測器}

$\mathsf S^\sharp$ の\emph{両側イデアル}とは，$\mathsf S^\sharp\mathsf I\mathsf S^\sharp\subseteq\mathsf I$ を満たす集合 $\mathsf I$ である．
文脈 $(u,v)$ が $\mathsf I$ を\emph{強制する}とは $\widehat s(u)\,\mathsf S^\sharp\,\widehat s(v)\subseteq\mathsf I$ のことをいう：何を入れても結果の型は $\mathsf I$ に入る．

\begin{proposition}[イデアルによる制限]\label{prop:ideal}
$s$ を型付け，$\mathsf I$ を $\mathsf S^\sharp$ の両側イデアルとし，$(G,(C_A)_A)$ を $L$ の型なしの $k$-文脈の提示で，
初期記号はどの右辺にも現れず，初期でないすべての非終端記号のすべての文脈が $\mathsf I$ を強制するものとする．
$G'$ を，補題~\ref{lem:product} の細分化 $G^s$ から，$G$ の初期記号 $S$ と $\mu\notin\mathsf I$ についての非終端記号 $(S,\mu)$ をその生成規則とともに除き，$S$ が初期で $\mu\in\mathsf I$ である $(S,\mu)$ を初期記号としたものとする．
このとき $L(G')=L'=L\cap\widehat s^{-1}(\mathsf I)$ であり，$L'$ に相対的に，族 $(C_A\Lc\cap F_\mu)_{(A,\mu)}$ は $G'$ の SPP である．
もとの提示が強ければ，これも強い．
よって $L'$ は，同じ文脈，高々 $|N||\mathsf S^\sharp|$ 個の非終端記号，長さ $\rho$ 以下の右辺の $s$-型付き $k$-文脈の提示を持つ．
\end{proposition}
\begin{proof}
$L$ と $L'$ に相対的な作用素を $C\Lc$ と $C^{\lhd'}$ と書く．初期記号はどの右辺にも現れないので，除去はほかの生成規則に影響せず，補題~\ref{lem:product} により $L(G')=L'$．$G'$ の非終端記号は，$G$ で初期でない $A$ についての $(A,\mu)$ か，初期の $S$ と $\mu\in\mathsf I$ についての $(S,\mu)$ である．
(i) $(u,v)$ が $\mathsf I$ を強制するなら，$uxv\in L$ から $uxv\in L'$ が従う．よって初期でないすべての $B$ について $C_B^{\lhd'}=C_B\Lc$．
(ii) $\mu\in\mathsf I$，$x\in F_\mu$ なら，すべての $(u,v)$ について $\widehat s(uxv)\in\mathsf S^\sharp\mu\mathsf S^\sharp\subseteq\mathsf I$．
よって初期記号 $S$ について $C_S^{\lhd'}\cap F_\mu=C_S\Lc\cap F_\mu$．
$G'$ の生成規則 $(A,\mu)\to w_0(A_1,\mu_1)w_1\cdots(A_m,\mu_m)w_m$ について，右辺の成分の積は，(i)，もとの前不動点，型の整合性により
$w_0C_{A_1}\Lc w_1\cdots C_{A_m}\Lc w_m\cap F_\mu\subseteq C_A\Lc\cap F_\mu$ に入り，(i) または (ii) により $C_A\Lc\cap F_\mu=C_A^{\lhd'}\cap F_\mu$．
初期の成分 $C_S\Lc\cap F_\mu$（$\mu\in\mathsf I$）は，もとの提示の健全性により $L\cap\widehat s^{-1}(\mathsf I)=L'$ に含まれるので，この族は SPP である．
強い場合には，(i) と (ii) から $C_A^{\lhd'}\cap F_\mu=L(G,A)\cap F_\mu=L(G',(A,\mu))$．
\end{proof}

\begin{remark}[二つの例]\label{rem:ideal-examples}
命題~\ref{prop:deletion-upper}(b) は，無用な非終端記号を除けば，$L=L_n$，$S_0$ とその生成規則を除いて $E_n$ を初期記号とした $G=G_n$，$s=s_U$，$\mathsf I=\{\zeta\}$ の場合である：
$i<n$ の各 $E_i$ の文脈 $(\Lce^{n-i},\Rce^{n-i})$ は標識を含むので $\{\zeta\}$ を強制する．
別の観測器として，飽和する加法の $\{0,1,2\}$ への $t(w)=\min(\#_{\Lce}(w),2)$ と $\mathsf I=\{2\}$ をとる．
$i\le n-2$ には文脈 $(\Lce^{n-i},\Rce^{n-i})$，$E_{n-1}$ には $(\Lfl,u_{n-1}\Rfl)$（$n\ge3$）を用いると，これらは標識を二つ以上含み $G_n$ の強 $1$-FCP を与えるので，
標識を二つ以上含む構造の言語は，長さ $2n+1$ 以下の文脈による $t$-型付き強 $1$-FCP の提示を持つ．
\end{remark}

\subsection*{削除族上の標識の観測器}

\begin{proposition}[標識の観測器のもとでの全長]\label{prop:mark-total}
$n\ge2$ とする．$H_n$ の $s_U$-型付きの枠の提示（命題~\ref{prop:deletion-upper}(b)）の全長は $O(n^2)$ である．
一方，$M_n$ のどの見本の提示も，型付きでもそうでなくても，全長 $2\sqrt{3\cdot2^n-n-5}$ 以上を持つ．
\end{proposition}
\begin{proof}
$H_n$ の提示は，右辺の長さが $4$ 以下の生成規則 $6n-4$ 個と，長さ $2n$ 以下の文脈 $2n$ 個を持つ．
二つ目の主張については $r_n=3\cdot2^n-n-5=|R|$ とする．
定理~\ref{thm:general-lower} の証明で用いた $z$ の導出木（$R$，$T$，$T^\circ$ は補題~\ref{lem:one-word} のもの）では，鎖の節点は互いに異なる非終端記号を持ち，したがって互いに異なる生成規則を用いる．
それらの右辺の長さの和を $q$ とする．その証明の被覆の論法は，$R$ の各断片をこれらの右辺のどれかの記号に割り当てるので，$r_n\le q\max(\kappa,1)$．
全長 $\Lambda$ はこれらの生成規則を数えて $q$ 以上となり，それとは別に最大長の核の語を数えて $\max(\kappa,1)$ 以上となる．
有限言語 $M_n$ では $\emptyset\RL=\emptyset$ なので，ある核は空でない．
よって $\Lambda\ge q+\max(\kappa,1)\ge2\sqrt{q\max(\kappa,1)}\ge2\sqrt{r_n}$．
\end{proof}

定理~C はこれらの結果をまとめたものである：(i) は命題~\ref{prop:deletion-upper}(b) と命題~\ref{prop:untyped-frames}，(ii) は定理~\ref{thm:deletion} と注意~\ref{rem:Ln-sU}，(iii) は命題~\ref{prop:mark-total}，(iv) は下の定理~\ref{thm:frame-total} である．

\begin{theorem}[一般の隔たり]\label{thm:general-gap}
$L$ を有限言語，$s$ を型付け，$\mathsf I$ を $\mathsf S^\sharp$ の両側イデアルとする．
$L$ が命題~\ref{prop:ideal} のような長さ $\ell$ 以下の文脈による型なしの $k$-文脈の提示を持ち，
$L\setminus\widehat s^{-1}(\mathsf I)=\{w\}$ であり，$w$，$R$，$T$，$T^\circ$ が補題~\ref{lem:one-word} の条件を満たすとする．
このとき $L\setminus\{w\}$ は，長さ $\ell$ 以下の文脈，長さ $\rho$ 以下の右辺，高々 $|N||\mathsf S^\sharp|$ 個の非終端記号の $s$-型付き $k$-文脈の提示を持つ．
一方，$L\setminus\{w\}$ のどの見本の提示も，型付きでもそうでなくても，$\rho|N|\max(\kappa,1)\ge|R|$ を満たし，全長 $2\sqrt{|R|}$ 以上を持つ．
\end{theorem}
\begin{proof}
$L\cap\widehat s^{-1}(\mathsf I)=L\setminus\{w\}$ なので，最初の主張は命題~\ref{prop:ideal} である．
二つ目の主張は補題~\ref{lem:one-word} と定理~\ref{thm:general-lower}，および命題~\ref{prop:mark-total} の証明による．$T=\eps$ なら，定理~\ref{thm:general-lower} の証明から $\kappa\ge|R|$ なので，全長は $|R|+1\ge2\sqrt{|R|}$ 以上である．
\end{proof}
定理~\ref{thm:deletion} と命題~\ref{prop:mark-total} の下界は，$L=L_n$，$s=s_U$，$\mathsf I=\{\zeta\}$，$w=w^*_n$，$\ell=2n$，$|R|=3\cdot2^n-n-5$ の場合である．

\subsection*{枠が必要とする観測器}
全長については，観測器は助けになるだけでなく必要でもある．
$1\le j<n$ について $v_j=\Lce w^*_{j-1}\Rce$ とおく．$v_j\in L(G_n,E_j)$ であり，$|w^*_j|-|v_j|=|w^*_{j-1}|$．
$w^*_n$ の $E_j$-節点は根からの道 $\pi\in\{l,r\}^{n-j}$ で区別される．節点 $\pi$ について $w^*_n=P_\pi w^*_jQ_\pi$ と書き，
$z_\pi=P_\pi v_jQ_\pi$ とおく．これは標識をちょうど一つ持つ $M_n$ の語である．

\begin{lemma}[標識付きの節点]\label{lem:marked-node}
$1\le j<n$，$\pi\in\{l,r\}^{n-j}$ とする．$G$ を SPP $(C_A\Lc)_A$ を持つ $M_n$ の CFG，
または $\widehat s(v_j)=\widehat s(w^*_j)$ となる $s$ について SPP $(C_A\Lc\cap F_{\tau(A)})_A$ を持つ $M_n$ の $s$-型付き CFG とし，
文脈の長さは $\ell$ 以下とする．このとき，$z_\pi$ のどの導出木のどの節点も，その収量 $x$ が $v_j$ の出現を含むならば，
$|x|\ge|w^*_n|-\ell-|w^*_{j-1}|$ を満たす．同じことだが，その終端の文脈 $(p,q)$ は $|p|+|q|\le\ell$ を満たす．
\end{lemma}
\begin{proof}
節点のラベルを $A$ とし，$x'$ を $x$ の中の $v_j$ の出現を $w^*_j$ で置き換えたものとする．このとき $px'q=w^*_n$ であり，
型付きの場合は $\widehat s(x')=\widehat s(x)$．もし $x'$ が $A$ の成分に入るなら，補題~\ref{lem:reach} から $w^*_n\in M_n$ が得られてしまう．
$x$ はその成分に入り，$x'$ は同じ型を持つので，ある $c\in C_A$ は $x$ を受け入れるが $x'$ を拒否する．
さて $c[x]\in M_n\subseteq L_n$ であり，$c[x]$ における $v_j$ の出現は $E_j$-節点である（補題~\ref{lem:nodes}）ので，
$c[x']\in L_n\setminus M_n$，すなわち $c[x']=w^*_n$．$|x'|=|x|+|w^*_{j-1}|$ なので
$|x|=|w^*_n|-|c|-|w^*_{j-1}|\ge|w^*_n|-\ell-|w^*_{j-1}|$ であり，$|z_\pi|=|w^*_n|-|w^*_{j-1}|$ から二つ目の形が得られる．
\end{proof}

\begin{lemma}[混合]\label{lem:mixing}
$1\le j<n$，$\pi,\pi'\in\{l,r\}^{n-j}$ とする．$P_\pi v_jQ_{\pi'}\in L_n$ となるのは $\pi=\pi'$ のときに限る．
\end{lemma}
\begin{proof}
$\pi=\pi'$ ならこの語は $z_\pi$ である．$P_\pi$ の閉じていない括弧は，節点 $\pi$ の祖先の $n-j$ 個の開き括弧 $\Lfl$ である．
深さ $k-1$ のものと次の閉じていない括弧（または $v_j$）のあいだには，$\pi_k=r$ なら道の左の兄弟（ある $L(G_n,E_i)$ の語）があり，$\pi_k=l$ なら何もない．
同様に $Q_{\pi'}$ では，右から $k$ 番目の閉じていない $\Rfl$ の前に，$\pi'_k=l$ なら右の兄弟がある．
$P_\pi v_jQ_{\pi'}\in L_n$ なら，これらの閉じていない括弧は入れ子の順に対応し，対応する対はそれぞれ導出木の節点である．
深さ $k-1$ の節点は，内側の節点（または $v_j$）と，現れている $P_\pi$ の左の兄弟と $Q_{\pi'}$ の右の兄弟を子に持つ．
$\Lfl$ の節点はちょうど二つの子を持つので，$\pi_k=r$ と $\pi'_k=l$ のちょうど一方が成り立ち，すなわち $\pi_k=\pi'_k$．
\end{proof}

\begin{theorem}[全長では枠は観測器を必要とする]\label{thm:frame-total}
$n\ge2$，$1\le j<n$ とし，$s$ を $\widehat s(v_j)=\widehat s(w^*_j)$ となる型付け，または型付けなしとする．
$M_n$ の（$s$-型付きの）どの枠の提示も，その文脈の長さ $|u|+|v|$ の最大値を $\ell$（文脈がなければ $\ell=0$）として
\[
(3\cdot2^{j-1}-1)\,(\ell+1)\sum_{A\to\alpha}|\alpha|\ \ge\ 2^{n-j}
\]
を満たす．特に，$M_n$ のどの型なしの枠の提示も，また $\widehat s(\Lce\ta\Rce)=\widehat s(\Lfl\ta\ta\Rfl)$ となる観測器のもとでのどの型付きの枠の提示も，
$(\ell+1)\sum_{A\to\alpha}|\alpha|\ge2^{n-2}$ を満たし，全長 $2^{n/2}$ 以上を持つ．
\end{theorem}
\begin{proof}
各 $\pi\in\{l,r\}^{n-j}$ について $z_\pi$ の導出木を一つ定め，収量が $v_j$ の出現全体を含む最も低い節点を $\nu_\pi$ とする．節点の子の収量は互いに交わらない区間なので，そのうち高々一つしかこの出現を含まず，$\nu_\pi$ は一意に定まる．
その生成規則を $A\to\sigma_1\cdots\sigma_r$，終端の文脈を $(p_\pi,q_\pi)$ とする．
補題~\ref{lem:marked-node} により $|q_\pi|\le\ell$ であり，$q_\pi$ は $z_\pi$ で $v_j$ の後にあるので，その長さの $w^*_n$ の接尾辞である．よって $|q_\pi|=|q_{\pi'}|$ なら $q_\pi=q_{\pi'}$．
$v_j$ の最初の記号 $\Lce$ を収量に含む記号を $\sigma_e$ とする．$\sigma_e$ が終端記号ならその収量はその記号であり，
そうでなければ $\nu_\pi$ の選び方からその収量は $v_j$ 全体を含まない．よって $\sigma_1\cdots\sigma_e$ の収量は，ある $1\le d\le|v_j|-1=3\cdot2^{j-1}-1$ について $v_j$ の最初の $d$ 記号の後で終わる．

$\pi$ と $\pi'$ が同じ生成規則，$e$，$d$，$|q|$ を与えるとする．
$z_\pi$ の木で，$\nu_\pi$ の下の $\sigma_{e+1},\dots,\sigma_r$ の部分木を $\nu_{\pi'}$ の下のものに替える．
これは $G$ の（型付きの場合は，ラベルが変わらないので $(G,\tau)$ の）導出木であり，
二つの木はどちらも $v_j$ を $d$ 記号の後で分け，$q_\pi=q_{\pi'}$ なので，その収量は $P_\pi v_jQ_{\pi'}$ である．
よって $P_\pi v_jQ_{\pi'}\in M_n\subseteq L_n$ であり，補題~\ref{lem:mixing} により $\pi=\pi'$．
生成規則と位置 $e$ の組は $\sum|\alpha|$ 個なので，$2^{n-j}\le\sum|\alpha|\cdot(3\cdot2^{j-1}-1)\cdot(\ell+1)$．

$j=1$ ではこれは $(\ell+1)\sum|\alpha|\ge2^{n-2}$ である．全長 $\Lambda$ は $\sum|\alpha|+\ell+1$ 以上である：
ある $C_A$ が空でなければ最長の文脈が $\ell+1$ を与え，そうでなければ $\ell=0$ で生成規則が $\sum|\alpha|+1$ 以上を与える．
よって $\Lambda\ge2\sqrt{(\ell+1)\sum|\alpha|}\ge2^{n/2}$．
\end{proof}

枠を用いるのは補題~\ref{lem:marked-node} だけである：見本の場合や $s_U$ のもとでの枠の場合には，$H_n$ のように節点 $\nu_\pi$ は木の深いところにあってよく，$|q_\pi|$ は抑えられない．
したがって，枠が多項式の全長を持つためには，観測器は $n$ が大きくなるとき固定したどの $j$ についても $v_j$ と $w^*_j$ を分けなければならない．これは必要条件であって十分条件ではない．
長さの偶奇の観測器は $\Lce\ta\Rce$ と $\Lfl\ta\ta\Rfl$ を分けるが，$\Lce w^*_1\Rce$ と $w^*_2$（長さ $6$ と $10$）は分けないので，
$n\ge3$ では $j=2$ の場合により，その $M_n$ の枠の提示もやはり全長 $2^{n/2}/\sqrt5$ 以上を持つ．
$j=1$ の場合が当てはまる提示のうち，枠の長さが $2n$ 以下のものの中では $\sum|\alpha|\ge2^{n-2}/(2n+1)$ なので，この類の中で $F_n$（$\sum|\alpha|=\Theta(n2^n)$）は $O(n^2)$ 倍の範囲で最適である．

因子 $\ell+1$ は除けない：命題~\ref{prop:deletion-upper}(c) の証明の $c_i$，$(p_i,q_i)$ について，型なしの $H_n$ に文脈 $\{c_i\}$（$(E_i,1)$），$\{c_i,(p_i,q_i)\}$（$(E_i,\zeta)$，$i<n$），$\{(\eps,\eps)\}$（$(E_n,\zeta)$）を与えたものは $M_n$ の $2$-文脈の提示である（補題~\ref{lem:nodes} により $p_ixq_i\in M_n$ は $x$ が標識を持つことと同値で，$\zeta$ の生成規則は標識を加えるか型 $\zeta$ の子を持つ）．このとき $\ell=3\cdot2^n-6$ だが，$n\ge10$ で $\sum|\alpha|\le4(6n-4)<2^{n-2}$．補題~\ref{lem:marked-node} の証明が枠を使うのは $x$ と $x'$ を区別する文脈を通してだけなので，成分が $|c|\le\ell$ なる集合 $\{c\}\Lc$ のブール結合であっても定理~\ref{thm:frame-total} はそのまま成り立つ（補題~\ref{lem:nodes} により，区別する $c$ は $x$ を受理する）．拡張正則の演算のもとでは成り立たず，$H_n$ は短い原子しか要しない（第~\ref{sec:applications} 節）．

\begin{example}[平らな例]\label{ex:flat}
一般的な主張も定理~\ref{thm:frame-total} の方法も木を必要としない．$\Sigma=\{a,b,\#\}$，$B_i=\{a,b\}^{2^i}$，$t_i=\#a^i\#$，$W_n=a^{2^n}$，
$L^\flat_n=B_n\cup\bigcup_{i=0}^nt_iB_i$，$M^\flat_n=L^\flat_n\setminus\{W_n\}$ とする．
\begin{enumerate}[label=\textup{(\roman*)},nosep]
\item $L^\flat_n$ の文法 $A_0\to a\mid b$，$A_i\to A_{i-1}A_{i-1}$，$S\to A_n\mid t_iA_i$ $(0\le i\le n)$ は，
$t_i$ が接頭辞符号をなすので，$A_i$ に文脈 $(t_i,\eps)$，$S$ に $(\eps,\eps)$ を与えて強 $1$-FCP を持つ（タグは \citealp[Thm.~5]{ky2021} の接尾辞 $\#d^i$ と同じ役割を果たす）．
$s^\flat$ を $b$ または $\#$ が現れるかどうかを記録する観測器とする．各 $t_i$ は $\#$ を含むのでその文脈は $\{\zeta\}$ を強制し，
$W_n$ は型 $1$ を持つ $L^\flat_n$ の唯一の語である．命題~\ref{prop:ideal} により，$M^\flat_n$ は全長 $O(n^2)$ の $s^\flat$-型付き強 $1$-FCP の提示を持つ．
\item $R=a^{2^n-1}$，$T=b$，$T^\circ=a$ とすると補題~\ref{lem:one-word} が適用できる：$yb\in L^\flat_n$ なら最後の $b$ は中身の語にあり，
それを $a$ に替えると $ya\in L^\flat_n$ となる．よって $M^\flat_n$ のどの見本の提示も，型なしでもどの観測器のもとでも，
$\rho|N|\max(\kappa,1)\ge2^n-1$ を満たし，全長 $2\sqrt{2^n-1}$ 以上を持つ（定理~\ref{thm:general-gap}）．
\item 観測器がないとき，または $\widehat s(a)=\widehat s(b)$ となる観測器のもとでは，$M^\flat_n$ のどの枠の提示も
$(\ell+1)\sum|\alpha|\ge2^n$ を満たし，全長 $2^{n/2+1}$ 以上を持つ．$z_i=a^iba^{2^n-1-i}$（$0\le i<2^n$）をとる．補題~\ref{lem:marked-node} と同様に，収量 $x$ がこの $b$ を含む節点は，$x$ を受け入れ $x'$（$b$ を $a$ に替えたもの）を拒否する文脈 $c$ を持つ．
$c[x]$ の $b$ は中身の語にあるので $c[x']=W_n$ であり，収量の外には高々 $\ell$ 個の記号しかない．
そのような最も低い節点では $b$ は終端記号である．$z_i$ と $z_{i'}$ が生成規則，$b$ の位置，右の外側の文脈の長さを共有するなら，
$b$ より右の部分木を交換すると $a^iba^{2^n-1-i'}$ が導出され，これが $M^\flat_n$ に入るのは $i=i'$ のときに限る．
\end{enumerate}
\end{example}

%======================================================================
\section{未解決問題}\label{sec:open}

\begin{openproblem}[弱 $1$-FCP の単調性]\label{op:weak-fcp}
$s\preceq h$ について $\tFCP1{w}{s}\subseteq\tFCP1{w}{h}$ か．
命題~\ref{prop:weak-not-preserved} は，すべての $k\ge2$ についてこの問いに否定的に答える．$k=1$ では有限言語は反例を与えない．
\end{openproblem}

\begin{openproblem}[頑健な分離]\label{op:robust}
大きさ 2 以上のアルファベット上のすべての型付け $s$ について $\tFCP{}{vw}{s}\ne\tFCP{}{w}{s}$ か．
\citet[Thm.~14]{ky2017} の証人 $L'$ については，$\$$ の指示型付けが収量 $a^{m_6}\$a^{l_6}$ を，\citet[Lemma~13]{ky2017} の証明で用いられる語 $c^id^ie^i$ から分けるので，
その証明はもはや適用できない．そのとき $L'$ が型付き弱 FCP を持つかどうかは分からない．
\end{openproblem}

\begin{openproblem}[観測器なしの枠の指数]\label{op:frame-exponent}
$M_n$ の型なしの枠の提示の最小の全長を $f(n)$ とする．定理~\ref{thm:frame-total} と命題~\ref{prop:untyped-frames} により $2^{n/2}\le f(n)=O(n2^n)$．
$\lim_n\log_2f(n)/n$ は存在するか．$\liminf_n\log_2f(n)/n$ と $\limsup_n\log_2f(n)/n$ はいくつか．特に，観測器なしの $M_n$ の枠と見本の提示は，全長で多項式倍を超えて異なるか．
一般の一語の削除に対する形には補題~\ref{lem:mixing} に代わるものが要る（例~\ref{ex:flat}(iii) は代わりに長さを用いる）．
\end{openproblem}

\begin{openproblem}[強 FKP]\label{op:further}
型付き強 FKP は，極弱 FKP と弱 FKP と同じく（命題~\ref{prop:duality}），細分化について反単調か．
\end{openproblem}

%======================================================================
\section{結論}\label{sec:conclusion}
枠と見本は，同じ統語カテゴリーを名指すのに指数的に異なりうる．
大域的な認可条件について，短い枠が非終端記号が少なく右辺の短い文法のすべてのカテゴリーを名指すのに対し，
どの見本の提示でも $\rho$，$|N|$，$\kappa$ の積は指数的であり，これが $\mathsf{Rec}_s$ と主学習器のデータを下から抑え，双対学習器は別の論法で抑えられる．
全長では，枠は標識の観測器のもとで短くなるが，$\Lce\ta\Rce$ と $\Lfl\ta\ta\Rfl$ が同じ型を持つと短くならない．見本はどの観測器のもとでも指数的なままである．
観測器を細分化すると，強 FCP と極弱 FCP については枠に有利にしか働かず，弱 FKP と極弱 FKP については見本に不利にしか働かない．弱 FCP は単調でない．

例~\ref{ex:toy} と，義務性および主格要件との類比は，数学的な概念を例示するためのものであり，自然言語について経験的な主張はしない．結果は最悪の場合の分離であって，すべての学習器についての主張ではない．
義務性はふつう頂点性（第一強勢は多くても一つ）と組にされる \citep{hyman2006}．標識をちょうど一つ持つ言語は本稿では扱わない．
一般の下界は有限言語と捕獲の条件~\eqref{eq:capture} を必要とする．補題~\ref{lem:one-word} は，削除される語 $RT^\circ$ について，ある $RT\in L$（$T\ne T^\circ$）がすべての $y$ で $yT\in L\Rightarrow yT^\circ\in L$ を満たすときにこれを与える．標識が二つ未満の構造をすべて削るときのように削除される語が複数あるとき，
枠の側は移る（注意~\ref{rem:ideal-examples}）が，見本の下界は得られていない．
不連続な構成素（多重文脈自由文法）と交差依存は，本稿の文脈自由の設定の範囲外である．\citet{kanazawa1998} の古典的な範疇文法は弱生成力では文脈自由だが，カテゴリーを型で名指すので，本稿では扱わない．

\section*{申告（Statements and Declarations）}

\paragraph{資金.}
本研究の実施にあたり資金の提供は受けていない．

\paragraph{競合する利益.}
著者に申告すべき競合する利益はない．

\paragraph{データの利用可能性.}
本稿で述べた研究には，実証的なデータを用いていない．

\paragraph{生成AIおよびAI支援技術の利用.}
本稿の準備と改訂にあたり，著者は言語の編集，原稿の構成，数学的な議論の検討の補助として OpenAI の ChatGPT と Anthropic の Claude を用いた．
著者は AI の補助による内容を検討・編集し，本稿で示した結果と証明に全責任を負う．

\appendix
\section{命題~\ref{prop:weak-not-preserved} の証明}\label{app:weak}

\begin{proof}[命題~\ref{prop:weak-not-preserved} の証明]
$G^+$ を，\citet[Lemma~12]{ky2017} の文法 $G'$ に生成規則 $S_1\to@$ を加えたものとする．
$S_1$ は自身の右辺以外のどの右辺にも現れないので，$L(G^+,S_1)=L_1\cup L_1^@$ であり，他の非終端記号はどれも言語を保ち，$L(G^+)=L^+$．
特に断らない限り，閉包はすべて $L^+$ に相対的である．
$L^+$ の語についての次の三つの事実を用いる：$c$，$d$ または $e$ を含む語は $L_1\cup L_1^@$ に属さない．
$a$，$b$，$d$ または $e$ で終わる語は $L_4\cup L_5$ に属さない．$d$ を含む語は $L_1\cup L_1^@\cup L_4$ に属さない．

$G'$ の非終端記号と，\citet[proofs of Lemmas~8 and~12]{ky2017} がそれらについて挙げる文脈集合 $C_A$ を表~\ref{tab:weak-spp} に示す．
ここで $\mathcal B=\{w\in V_0:|w|_a=|w|_b\}$ とする．第 3 列は $L^+$ に相対的な $C_A\Lc$ であり，以下で示す
（長さ $5$ 以下のすべての語について機械でも確かめてある）．
\begin{table}[tb]
\centering\footnotesize
\caption{命題~\ref{prop:weak-not-preserved} の証明における $G'$ の非終端記号．}\label{tab:weak-spp}
\begin{tabular}{@{}llll@{}}
\toprule
$A$ & $C_A$ & $C_A\Lc$ & 特徴づけ\\
\midrule
$Q$ & $(\eps,\#cd)$, $(a,b\#de)$ & $\mathcal B$ & 強\\
$F$ & $(\eps,cd)$, $(\eps,ade)$ & $\mathcal B\#V$ & 強\\
$H$ & $(a\#c,d)$ & $\{c^id^i:i\ge0\}$ & 強\\
$E$ & $(a\#cd,e)$ & $e^*$ & 強\\
$C$ & $(b\#c,de)$ & $c^*$ & 強\\
$J$ & $(b\#d,e)$ & $\{d^ie^i:i\ge0\}$ & 強\\
$T$ & $(\$,a\%)$ & $V_0$ & 強\\
$U$ & $(e,c\%)$ & $\{e^jzc^k:z\in V_0,\ j\ne k\}$ & 強\\
$S_{4,5}$ & $(\eps,f)$ & $L_4\cup L_5$ & 強\\
$S_2$ & $(\eps,\eps)$, $(a\#,b)$ & $L_2\cup J_0$ & 弱\\
$S_3$ & $(\eps,\eps)$, $(b\#,a)$ & $L_3\cup J_0$ & 弱\\
$S_1$ & $(a\#,a)$, $(b\#,b)$ & $L_1\cup L_1^@\cup J_0$ & 弱\\
\bottomrule
\end{tabular}
\end{table}

\emph{主張.}  $Y\in\{L_2,L_3\}$ について，$u(Y\setminus\{\eps\})v\subseteq L^+$ となるどの文脈 $(u,v)$ も
$u\in V$，$v\in V_0$，$uv\in L^+$ を満たす．
特に $(Y\setminus\{\eps\})\R=Y\R$ であり，すべての $(u,v)\in Y\R$ について $uJ_0v\subseteq J_0$．

$Y=L_2$ については語 $x_t=\#cde^ta\in L_2$（$t\ge0$）を用いる．
各 $ux_tv$ は $d$ を含むので $L_2\cup L_3\cup L_5$ に属する．
もし $v$ が $\%$ を含むなら，すべての $ux_tv$ は $L_5$ に属する．そこでは $x_t$ の $\#$ は $V$ に属する接頭辞の部分にあるので，
$cd$ は $i=1$ のブロック $c^id^i$，$e^t$ は $j=t$ のブロック $e^j$ であり，$v=z'c^k\%f^l$ で $k$ は $v$ によって決まる．
$t=k$ とすると $j\ne k$ に反する．
$u$ に $\%$ があることは不可能である．$\%$ の後には $f$ しか来ないからである．
よって $ux_tv\in L_2\cup L_3\subseteq V\,c^*d^*e^*V_0$ であり，$x_t$ は $\#$ で始まり $a$ で終わるので，$u\in V$ かつ $v\in V_0$．
$Y=L_3$ については $x=\#cd^2e^2a\in L_3$ を用いる：
$L_5$ では $\#$ の後の $c$ と二つの $d$ がブロック $c^id^i$ をなさなければならないが，それは不可能なので，同じ結論が従う．
すると $uv\in V\subseteq L_2\cup L_3$ なので，$(u,v)$ は $\eps$ も受け入れる．反単調性からもう一方の包含が得られる．
最後に $u(v'c^id^ie^iz')v=(uv')c^id^ie^i(z'v)\in J_0$．

\emph{強く特徴づけられる非終端記号.}  $A\in\{Q,F,H,E,C,J\}$（それぞれ $A\in\{T,U,S_{4,5}\}$）について，
\citet[proofs of Lemmas~8 and~12]{ky2017} は，$L_{123}$（それぞれ $L'$）に相対的に $|C_A|\le2$ なる $C_A$ で $L(G^+,A)=C_A\Lc$ を与える．
$Q,F,H,E,C,J$ について挙げられた各文脈は $c$，$d$ または $e$ を含み，右成分が $d$ または $e$ で終わるので，その $L^+$ での完成は $L_{123}$ に属する．
$T$，$U$，$S_{4,5}$ について挙げられた各文脈は $\%$ を含むか $f$ で終わるので，その $L^+$ での完成は $L_4\cup L_5$ に属する．
よって $L^+$ に相対的に $L(G^+,A)=C_A\Lc$ である：これらの非終端記号は強く特徴づけられる．

\emph{$S_2$ と $S_3$.}  $C_2=\{(\eps,\eps),(a\#,b)\}$ とおく．
完成 $a\#xb\in L^+$ は $b$ で終わり最初のブロックが $a$ なので，$L_4\cup L_5$ にも $L_1\cup L_1^@$ にも属さない．
よって $a\#xb\in L_2\cup L_3$ であり，$x$ は $\$,\%,f,@$ のどれも含まず，$x\in L^+$ から $x\in L_{123}$ が得られる．
したがって $L^+$ に相対的な $C_2\Lc$ は $L_{123}$ に相対的な $C_2\Lc$ に等しく，それは $L_2\cup J_0$ である \citep[proof of Lemma~8]{ky2017}．
$L_2\subseteq C_2\Lc$ なので $L_2\RL\subseteq C_2\Lc$．逆に $L_2\subseteq L_2\RL$ であり，主張により $J_0\subseteq L_2\RL$．
よって $L(G^+,S_2)\RL=C_2\Lc$ であり，対称的に $L(G^+,S_3)\RL=\{(\eps,\eps),(b\#,a)\}\Lc=L_3\cup J_0$．

\emph{$S_1$.}  $M$ を鏡映の文脈 $(w_1\#\cdots\#w_n,w_n^R\cdots w_1^R)$，$n\ge1$，$w_i\in V_0$ の集合とする．
$@$ を含む $L^+$ の語はちょうど $L_1^@$ の語なので，$L^+$ に相対的な $L_1^@$ の文脈は，$L_1^@$ に相対的な $L_1^@$ の文脈である．
$@$ を $\$$ に名前を替え，$L_{123}$ に相対的に $L(G,S_1)\R=M$ であること \citep[proof of Lemma~8]{ky2017} を用いると，この集合は $M$ である．
$L'$ に相対的には，$L_1$ の文脈は $M$ と文脈 $(v,zc^k\%f^l)$ からなり \citep[proof of Lemma~12]{ky2017}，
後者は $L_1^@$ の文脈ではない．よって $L^+$ に相対的に $L(G^+,S_1)\R=M$．
$@$ を含まない語については，$L^+$ に相対的な閉包 $M\Lc$ は $L'$ に相対的なものと一致し，
それは $L_1\cup J_0=\{(a\#,a),(b\#,b)\}\Lc$ である \citep[proof of Lemma~12]{ky2017}．
関わる文脈が $@$ を含まないからである．
$@$ を含む語 $x$ が $M\Lc$ に属するのは $x\in L_1^@$ のとき，かつそのときに限り，
$\{(a\#,a),(b\#,b)\}\Lc$ に属するのは $a\#xa,b\#xb\in L_1^@$ のとき，すなわちやはり $x\in L_1^@$ のとき，かつそのときに限る．
よって $L(G^+,S_1)\RL=L_1\cup L_1^@\cup J_0=\{(a\#,a),(b\#,b)\}\Lc$．これで $L^+\in\FCP2{w}$ が示された．

\emph{$s_0$-型付き文法.}  $G^+$ を $s_0$ について補題~\ref{lem:product} で細分化し，有用な非終端記号を残す．
非終端記号 $(A,\one)$ は $\{\eps\}=\emptyset\Lc\cap F_{\one}$ を生成する．
非終端記号 $(A,0)$ は $Y_A\setminus\{\eps\}$，$Y_A=L(G^+,A)$ を生成する．
$A$ が強く特徴づけられるなら，$C_A\subseteq Y_A\R\subseteq(Y_A\setminus\{\eps\})\R$ と $Y_A\setminus\{\eps\}=C_A\Lc\cap\Sigma^+$ から
$\cl_0(Y_A\setminus\{\eps\})=C_A\Lc\cap\Sigma^+$ が得られる．
非終端記号 $S_1$ は空語を生成しないので，$\cl_0(Y_{S_1})=L(G^+,S_1)\RL\cap\Sigma^+$．
$S_2,S_3$ については，主張から $(Y_A\setminus\{\eps\})\RL=Y_A\RL=C_A\Lc$ が得られるので，$\cl_0(Y_A\setminus\{\eps\})=C_A\Lc\cap\Sigma^+$．
よって細分化した文法のどの非終端記号も高々 2 個の文脈で特徴づけられ，$L^+\in\tFCP2{w}{s_0}$．

(b) $h$-型付きの $(G,\tau)$ が $L^+$ を生成し，型付き弱 FCP を持つと仮定する．
$\zeta$ は吸収元なので，左辺の型が $1$ または $\one$ の生成規則は型 $\zeta$ の子を持たない．
型 $\zeta$ の非終端記号と，その中の初期記号を除くと，$L(G_0)=L^+\cap\Sigma_0^*=L'$ となる CFG $G_0$ が残り，
残ったどの非終端記号も言語を保つ．
\citet[Lemma~13]{ky2017} の証明は $L'$ の任意の CFG に適用でき，Ogden の補題 \citep{ogden1968} を用いる．
これを $G_0$ に適用すると，収量 $a^{m}\$a^{l}$ を持つ（型 $1$ の）非終端記号 $B$ と数 $k_1$ が得られ，
$i>k_1$ について $L(G_0,B)\RL$ は $Vc^id^ie^iV_0$ の語を含まない．
一方，有限のどの $W\subseteq L(G_0,B)\R$ についても，$i>k_2$ なるすべての $i$ について $c^id^ie^i\in W\Lc$ となる $k_2$ がある
（いずれも $L'$ に相対的）．
よって有限のどの $W\subseteq L(G_0,B)\R$ も $L(G_0,B)\RL\cap\Sigma_0^+=W\Lc\cap\Sigma_0^+$ を満たさない．
$L^+$ に相対的に，$L(G,B)$ のどの文脈も $\$$ を含む語を受け入れるので，$@$ を含まない．
$@$ を含まない語の上では二つの言語は一致する．
よって $L^+$ に相対的な $\cl_1(L(G,B))$ は，$L'$ に相対的な $L(G_0,B)\RL\cap\Sigma_0^+$ である．
$C$ が有限の型付き弱特徴づけ $\cl_1(L(G,B))=C\Lc\cap F_1$ があれば $L(G,B)\subseteq C\Lc$，すなわち $L^+$ に相対的に $C\subseteq L(G,B)\R$ である．
よって $C$ は $@$ を含まず，$L'$ に相対的に $C\subseteq L(G_0,B)\R$ であり，$W=C$ は $B$ の選び方に反する．

最後の主張は，(a)，(b) と，$k\ge2$ について $\tFCP2{w}{s_0}\subseteq\tFCP{k}{w}{s_0}$ であることから従う．
\end{proof}

\begin{thebibliography}{99}

\bibitem[Bod(2006)]{bod2006}
R.~Bod.
Exemplar-based syntax: How to get productivity from examples.
\textit{The Linguistic Review} 23(3):291--320, 2006.

\bibitem[Bybee(2006)]{bybee2006}
J.~Bybee.
From usage to grammar: The mind's response to repetition.
\textit{Language} 82(4):711--733, 2006.

\bibitem[Clark(2010)]{clark2010icgi}
A.~Clark.
Learning context free grammars with the syntactic concept lattice.
In \textit{ICGI 2010}, LNAI 6339, pp.~38--51, Springer, 2010.

\bibitem[Clark(2011)]{clark2009fg}
A.~Clark.
A learnable representation for syntax using residuated lattices.
In \textit{Formal Grammar (FG 2009)}, LNAI 5591, pp.~183--198, Springer,
2011.

\bibitem[Clark(2015)]{clark2013}
A.~Clark.
The syntactic concept lattice: another algebraic theory of the
context-free languages?
\textit{Journal of Logic and Computation} 25(5):1203--1229, 2015
(online 2013).

\bibitem[Clark(2016)]{clark2016icgi}
A.~Clark.
Testing distributional properties of context-free grammars.
In \textit{ICGI 2016}, PMLR 57, pp.~42--53, 2016.

\bibitem[Clark and Eyraud(2007)]{clark-eyraud2007}
A.~Clark and R.~Eyraud.
Polynomial identification in the limit of substitutable context-free
languages.
\textit{Journal of Machine Learning Research} 8:1725--1745, 2007.

\bibitem[Clark and Yoshinaka(2016)]{clark-yoshinaka2016}
A.~Clark and R.~Yoshinaka.
Distributional learning of context-free and multiple context-free
grammars.
In \textit{Topics in Grammatical Inference}, pp.~143--172, Springer, 2016.

\bibitem[Coste and Nicolas(2019)]{coste-nicolas2019}
F.~Coste and J.~Nicolas.
Learning local substitutable context-free languages from positive examples
in polynomial time and data by reduction.
In \textit{ICGI 2019}, PMLR 93, pp.~155--168, 2019.

\bibitem[Fries(1952)]{fries1952}
C.~C.~Fries.
\textit{The Structure of English: An Introduction to the Construction of
English Sentences}.
Harcourt, Brace, New York, 1952.

\bibitem[Gold(1967)]{gold1967}
E.~M. Gold.
Language identification in the limit.
\textit{Information and Control}, 10(5):447--474, 1967.

\bibitem[Harris(1954)]{harris1954}
Z.~S.~Harris.
Distributional structure.
\textit{Word} 10(2--3):146--162, 1954.

\bibitem[Hyman(2006)]{hyman2006}
L.~M.~Hyman.
Word-prosodic typology.
\textit{Phonology} 23(2):225--257, 2006.

\bibitem[Kanazawa(1998)]{kanazawa1998}
M.~Kanazawa.
\textit{Learnable Classes of Categorial Grammars}.
CSLI Publications, Stanford, 1998.

\bibitem[Kanazawa(2023)]{kanazawa2023wollic}
M.~Kanazawa.
Learning context-free grammars from positive data and membership queries.
In \textit{WoLLIC 2023}, LNCS 13923, pp.~3--18, Springer, 2023.

\bibitem[Kanazawa and Yoshinaka(2017)]{ky2017}
M.~Kanazawa and R.~Yoshinaka.
The strong, weak, and very weak finite context and kernel properties.
In \textit{LATA 2017}, LNCS 10168, pp.~77--88, Springer, 2017.

\bibitem[Kanazawa and Yoshinaka(2021)]{ky2021}
M.~Kanazawa and R.~Yoshinaka.
A hierarchy of context-free languages learnable from positive data and
membership queries.
In \textit{ICGI 2021}, PMLR 153, pp.~18--31, 2021.

\bibitem[Kanazawa and Yoshinaka(2023)]{ky2023}
M.~Kanazawa and R.~Yoshinaka.
Extending distributional learning from positive data and membership
queries.
In \textit{ICGI 2023}, PMLR 217, pp.~8--22, 2023.

\bibitem[Kishimoto(2010)]{kishimoto2010}
H.~Kishimoto.
Subjects and constituent structure in Japanese.
\textit{Linguistics} 48(3):629--670, 2010.

\bibitem[Kuno(1973)]{kuno1973}
S.~Kuno.
\textit{The Structure of the Japanese Language}.
MIT Press, Cambridge, MA, 1973.

\bibitem[Kuriyama(2026)]{kuriyama-tcs1}
T.~Kuriyama.
Distributional learning of context-free languages under fixed
finite-semigroup typing.
Manuscript under review for \textit{Theoretical Computer Science}, 2026;
arXiv:1409.6247v5.

\bibitem[Lei{\ss}(2014)]{leiss2014}
H.~Lei{\ss}.
Learning context free grammars with the finite context property: a
correction of A.~Clark's algorithm.
In \textit{Formal Grammar 2014}, LNCS 8612, pp.~121--137, Springer, 2014.

\bibitem[Marcus(1993)]{marcus1993}
G.~F. Marcus.
Negative evidence in language acquisition.
\textit{Cognition}, 46(1):53--85, 1993.

\bibitem[Nespor and Vogel(1986)]{nespor-vogel1986}
M.~Nespor and I.~Vogel.
\textit{Prosodic Phonology}.
Foris, Dordrecht, 1986.

\bibitem[Ogden(1968)]{ogden1968}
W.~Ogden.
A helpful result for proving inherent ambiguity.
\textit{Math.\ Systems Theory} 2(3):191--194, 1968.

\bibitem[Selkirk(1984)]{selkirk1984}
E.~O.~Selkirk.
\textit{Phonology and Syntax: The Relation between Sound and Structure}.
MIT Press, 1984.

\bibitem[Yoshinaka(2011)]{yoshinaka2011dlt}
R.~Yoshinaka.
Towards dual approaches for learning context-free grammars based on
syntactic concept lattices.
In \textit{DLT 2011}, LNCS 6795, pp.~429--440, Springer, 2011.

\bibitem[Yoshinaka(2012)]{yoshinaka2012lata}
R.~Yoshinaka.
Integration of the dual approaches in the distributional learning of
context-free grammars.
In \textit{LATA 2012}, LNCS 7183, pp.~538--550, Springer, 2012.

\bibitem[Yoshinaka(2015)]{yoshinaka2015}
R.~Yoshinaka.
General perspective on distributionally learnable classes.
In \textit{MoL 2015}, pp.~87--98, 2015.

\end{thebibliography}
\end{document}
